\documentclass[11pt,a4paper]{article}
\usepackage[utf8]{inputenc}
\usepackage[margin=1in]{geometry}
\usepackage{amsmath,amssymb} % amssymb provides \checkmark
\usepackage{booktabs}
\usepackage{multirow}
\usepackage{hyperref}
\usepackage{xcolor}

% Custom commands
\newcommand{\Msun}{M_\odot}
\newcommand{\kms}{\,\mathrm{km\,s^{-1}}}
\newcommand{\pc}{\,\mathrm{pc}}
\newcommand{\cmcubed}{\,\mathrm{cm^{-3}}}

\title{Observational Constraints for IMBH-Cloud Interaction Simulations:\\
Summary of Oka et al.~(2017) and Ballone et al.~(2018)}
\author{Literature Review for SPH Simulation Setup}
\date{\today}

\begin{document}
\maketitle

\tableofcontents
\newpage

%===============================================================================
\section{Introduction}
%===============================================================================

This document summarizes the key observational constraints and simulation parameters from two papers studying the High-Velocity Compact Cloud (HVCC) CO--0.40--0.22, a candidate site for an intermediate-mass black hole (IMBH) in the Galactic Center region.

\begin{itemize}
    \item \textbf{Oka et al.~(2017)} -- Original ALMA observation and N-body simulation
    \item \textbf{Ballone et al.~(2018)} -- Analytical model and hydrodynamic simulation
\end{itemize}

The goal is to identify which parameters are \textcolor{blue}{observationally constrained} (MUST MATCH) versus \textcolor{red}{uncertain/model-dependent} (CAN VARY) for setting up SPH simulations.

%===============================================================================
\section{Oka et al.~(2017): Original Observation}
%===============================================================================

\subsection{Observational Data}

ALMA Band 6 observations of CO--0.40--0.22 at 230.5 GHz (CO $J$=2--1) and 265.9 GHz (HCN $J$=3--2).

\subsubsection{Cloud-Scale Properties}
\begin{itemize}
    \item Total cloud size: $\sim 5\pc$ (projected)
    \item Total velocity width: $\sim 100\kms$
    \item High CO $J$=3--2/$J$=1--0 ratio: $\geq 1.5$
    \item Projected distance from Sgr~A$^*$: $\sim 60\pc$
\end{itemize}

\subsubsection{Dense Clump Properties (ALMA resolution)}
\begin{itemize}
    \item Size: $\sim 0.3\pc$ (compact)
    \item Velocity width: $\sim 100\kms$
    \item Main body velocity range: $V_\mathrm{LSR} = -80$ to $-40\kms$
    \item High-velocity extensions: $V_\mathrm{LSR} = -105$ to $-5\kms$
    \item \textbf{Mass}: $M_\mathrm{clump} = 40\,\Msun$ (from HCN $J$=3--2 intensity)
    \item \textbf{Temperature}: $T_\mathrm{k} = 60\,\mathrm{K}$
    \item \textbf{Density}: $n(\mathrm{H}_2) = 10^{6.5}\cmcubed$
\end{itemize}

\subsubsection{Virial Analysis}
The clump is \textbf{NOT self-bound}:
\begin{align}
    M_\mathrm{actual} &= 40\,\Msun \\
    M_\mathrm{virial} &\simeq 4.1 \times 10^3\,\Msun \quad (\text{from } S = 0.15\pc, \sigma_V = 22\kms)
\end{align}
The factor of $\sim 100$ discrepancy indicates external forces (tidal interaction) dominate over self-gravity.

\subsubsection{Point Source CO--0.40--0.22$^*$}
\begin{itemize}
    \item Position: $(l, b) = (-0.3983^\circ, -0.2235^\circ)$
    \item Flux densities: $8.38 \pm 0.34$ mJy (231 GHz), $9.91 \pm 0.74$ mJy (266 GHz)
    \item Luminosity: $\nu L_\nu \sim 10^{32.3}\,\mathrm{erg\,s^{-1}}$ at $d = 8.3\,\mathrm{kpc}$
    \item Spectral index: $\alpha = 1.18 \pm 0.65$ (inverted, consistent with Sgr~A$^*$ at 1/500 luminosity)
    \item Size upper limit: $< 0.022\pc$ (beam minor axis)
\end{itemize}

\subsection{N-body Simulation Parameters}

Oka et al.~performed N-body simulations to reproduce the observed kinematics:

\begin{table}[h]
\centering
\caption{Oka et al.~(2017) N-body Simulation Setup}
\begin{tabular}{lll}
\toprule
\textbf{Parameter} & \textbf{Value} & \textbf{Notes} \\
\midrule
$M_\mathrm{BH}$ & $10^5\,\Msun$ & Point mass \\
$N_\mathrm{particles}$ & 1000 & Each $1\,\Msun$ \\
$\sigma_r$ (radial dispersion) & $0.2\pc$ & Gaussian distribution \\
$\sigma_v$ (velocity dispersion) & $1.43\kms$ & Initially virialized \\
$(X_0, Y_0)$ & $(9.8\pc, -0.65\pc)$ & Initial position \\
$(v_{X,0}, v_{Y,0})$ & $(-8.19\kms, 0.4\kms)$ & Initial velocity \\
Inclination $i$ & $70^\circ$ & Orbital plane to LOS \\
Position angle & $41.6^\circ$ & Of orbital plane \\
$V_\mathrm{LSR}$ (center of gravity) & $-120\kms$ & Line-of-sight velocity \\
Time step & $4.6 \times 10^3\,\mathrm{yr}$ & Leapfrog integration \\
Best-fit time & $t = 7.2 \times 10^5\,\mathrm{yr}$ & Post-pericenter \\
\bottomrule
\end{tabular}
\end{table}

\textbf{Key Result}: The simulation reproduced both the bulk kinematics of CO--0.40--0.22 and the compact dense clump near the continuum source.

%===============================================================================
\section{Ballone et al.~(2018): Analytical Model}
%===============================================================================

\subsection{Analytical Constraint on $M_\mathrm{BH}$}

Ballone et al.~derived a lower limit on the IMBH mass from the observed velocity gradient.

\subsubsection{Observed Velocity Gradient}
From the position-velocity diagram:
\begin{equation}
    \frac{\Delta v_\mathrm{LOS}}{\Delta r_\perp} \approx \frac{20\text{--}40\kms}{0.3\pc}
\end{equation}
where $\Delta v_\mathrm{LOS}$ is the line-of-sight velocity difference and $\Delta r_\perp$ is the projected separation.

\subsubsection{Tidal Acceleration Analysis}
For a cloud falling radially toward a point mass:
\begin{equation}
    \Delta v_\mathrm{LOS} = \frac{2 G M_\mathrm{BH}}{r^2} \cdot \Delta r \cdot \Delta t \cdot \cos^2\theta
\end{equation}

Under the most conservative assumption (radial infall, $\theta = 0$):
\begin{equation}
    M_\mathrm{BH} > \text{few} \times 10^4\,\Msun
\end{equation}

This is a \textbf{robust lower limit} independent of detailed simulation parameters.

\subsection{SPH Simulation Parameters}

Ballone et al.~performed 3D SPH simulations:

\begin{table}[h]
\centering
\caption{Ballone et al.~(2018) SPH Simulation Setup}
\begin{tabular}{lll}
\toprule
\textbf{Parameter} & \textbf{Value} & \textbf{Notes} \\
\midrule
$M_\mathrm{BH}$ & $5.7 \times 10^4\,\Msun$ & Point mass (minimum estimate) \\
$M_\mathrm{cloud}$ & $40\,\Msun$ & From HCN observations \\
$R_\mathrm{cloud}$ (initial) & $0.14\pc$ & Uniform density sphere \\
$T_\mathrm{cloud}$ & $60\,\mathrm{K}$ & Isothermal \\
$N_\mathrm{particles}$ & $10^5$ & SPH particles \\
Self-gravity & OFF & Tidal forces dominate \\
Initial velocity & Radial infall & Toward BH \\
Impact parameter & 0 & Head-on collision \\
\bottomrule
\end{tabular}
\end{table}

\textbf{Key Result}: The simulation reproduced the observed velocity gradient with $M_\mathrm{BH} = 5.7 \times 10^4\,\Msun$, representing the minimum mass required.

%===============================================================================
\section{Parameter Classification for SPH Simulations}
%===============================================================================

\subsection{MUST MATCH: Observationally Constrained Parameters}
\label{sec:must_match}

These parameters are directly measured from ALMA observations:

\begin{table}[h]
\centering
\caption{\textcolor{blue}{Observationally Constrained Parameters}}
\begin{tabular}{lll}
\toprule
\textbf{Parameter} & \textbf{Value} & \textbf{Source} \\
\midrule
Cloud mass & $M_\mathrm{cloud} = 40\,\Msun$ & HCN $J$=3--2 intensity \\
Cloud temperature & $T = 60\,\mathrm{K}$ & Molecular line ratios \\
Dense clump size & $\sim 0.3\pc$ & ALMA imaging \\
Velocity width & $\Delta v \sim 100\kms$ & HCN spectral line \\
Velocity gradient & $\sim 70\text{--}130\kms/\pc$ & Position-velocity diagram \\
Not self-bound & $M_\mathrm{vir}/M_\mathrm{actual} \sim 100$ & Virial analysis \\
\bottomrule
\end{tabular}
\end{table}

\textbf{Critical constraint}: The simulation output must produce a velocity gradient of $\Delta v_\mathrm{LOS} \approx 20$--$40\kms$ across a $\sim 0.3\pc$ region.

\subsection{CAN VARY: Model-Dependent Parameters}
\label{sec:can_vary}

These parameters are not directly constrained by observations:

\begin{table}[h]
\centering
\caption{\textcolor{red}{Model-Dependent Parameters}}
\begin{tabular}{lll}
\toprule
\textbf{Parameter} & \textbf{Range} & \textbf{Notes} \\
\midrule
$M_\mathrm{BH}$ & $10^4$--$10^5\,\Msun$ & Lower limit: few$\times 10^4\,\Msun$ \\
Initial cloud radius & $0.1$--$0.3\pc$ & Not directly observed \\
Impact parameter & 0--few pc & Trajectory unknown \\
Initial cloud shape & Sphere, Gaussian, etc. & Not constrained \\
Initial velocity & Variable & Depends on trajectory \\
Orbital inclination & $\sim 70^\circ$ & Fitted to reproduce PV diagram \\
\bottomrule
\end{tabular}
\end{table}

\subsection{Derived Constraints}

From the observational constraints, we can derive:

\subsubsection{Sound Speed}
For isothermal gas at $T = 60\,\mathrm{K}$ with $\mu = 2.33$ (molecular):
\begin{equation}
    c_s = \sqrt{\frac{k_B T}{\mu m_H}} \approx 0.46\kms
\end{equation}

\subsubsection{Initial Density}
For a uniform sphere with $M = 40\,\Msun$ and $R = 0.14\pc$:
\begin{equation}
    \rho_0 = \frac{3M}{4\pi R^3} \approx 3.5 \times 10^{-20}\,\mathrm{g\,cm^{-3}}
\end{equation}
\begin{equation}
    n(\mathrm{H}_2) \approx 10^4\cmcubed
\end{equation}

Note: The observed post-encounter density ($10^{6.5}\cmcubed$) is much higher, indicating significant compression.

\subsubsection{Tidal Radius}
At pericenter distance $r_p$, the tidal radius is:
\begin{equation}
    r_t = R_\mathrm{cloud} \left(\frac{M_\mathrm{cloud}}{3 M_\mathrm{BH}}\right)^{1/3}
\end{equation}

For $M_\mathrm{BH} = 10^5\,\Msun$, $M_\mathrm{cloud} = 40\,\Msun$:
\begin{equation}
    r_t \approx 0.07 \cdot R_\mathrm{cloud}
\end{equation}

The cloud is deeply in the tidal regime ($r_t \ll R_\mathrm{cloud}$).

%===============================================================================
\section{First-Principles Analysis of Initial Conditions}
%===============================================================================

This section analyzes the constraints on preparing a hydrostatic initial cloud
condition, examining compatibility with Koyama \& Inutsuka (2000) thermal
equilibrium and Bonnor-Ebert sphere physics.

\subsection{KI2000 Thermal Equilibrium: Fundamental Incompatibility}

The Koyama \& Inutsuka (2000) ISM cooling model provides equilibrium temperature
as a function of density. From the embedded data tables:

\begin{table}[h]
\centering
\caption{KI2000 Equilibrium Temperature vs Density (N$_H$ = 10$^{19}$ cm$^{-2}$)}
\begin{tabular}{lll}
\toprule
\textbf{Density $n$ [cm$^{-3}$]} & \textbf{$T_\mathrm{eq}$ [K]} & \textbf{Phase} \\
\midrule
0.01 & 697 & Warm Neutral Medium \\
0.2 & 562 & WNM \\
1 & 53--70 & \textcolor{red}{Unstable transition} \\
10 & 14 & Cold Neutral Medium \\
100 & 8 & CNM \\
1000 & 7.5 & CNM \\
$10^4$ & 4.4 & CNM \\
\bottomrule
\end{tabular}
\end{table}

\textbf{Critical Finding}: In KI2000 equilibrium:
\begin{itemize}
    \item $T = 60\,\mathrm{K}$ occurs \textit{only} at $n \sim 0.5$--$1\cmcubed$ (unstable region)
    \item At high density ($n > 10\cmcubed$): $T_\mathrm{eq} \approx 10$--$15\,\mathrm{K}$
    \item The observed $T = 60\,\mathrm{K}$ at high density is \textbf{NOT} KI2000 equilibrium
\end{itemize}

\textbf{Physical interpretation}: The observed 60~K temperature arises from
\textbf{shock heating during tidal compression}, as analyzed in detail in Section~\ref{sec:shock_heating}.

\subsection{Bonnor-Ebert Sphere Physics}

For an isothermal self-gravitating sphere in hydrostatic equilibrium:

\begin{align}
    \text{Sound speed:} \quad c_s &= \sqrt{\frac{k_B T}{\mu m_H}} \\
    \text{Scale length:} \quad \lambda &= \frac{c_s}{\sqrt{4\pi G \rho_c}} \\
    \text{Cloud radius:} \quad R &= \xi_s \cdot \lambda \\
    \text{Jeans mass:} \quad M_J &\sim \frac{c_s^3}{\sqrt{G^3 \rho_c}}
\end{align}

where $\xi_s$ is the dimensionless truncation radius ($\xi_s = 6.45$ for critical BE sphere).

\subsubsection{Constraint: Cannot Simultaneously Match M, R, and T}

For $T = 60\,\mathrm{K}$: $c_s = 0.46\kms$

\textbf{Case A: Fix $R = 0.14\pc$, $T = 60\,\mathrm{K}$} (Ballone configuration)
\begin{align}
    \lambda &= R/\xi_s = 0.14\pc / 5 = 0.028\pc \\
    \rho_c &= \frac{c_s^2}{4\pi G \lambda^2} \approx 3.4 \times 10^{-19}\,\mathrm{g\,cm^{-3}} \\
    n_c &\approx 8.7 \times 10^4\cmcubed \\
    M_\mathrm{BE} &\approx 1.5\text{--}3\,\Msun \quad \textcolor{red}{\neq 40\,\Msun}
\end{align}

\textbf{Case B: Fix $M = 40\,\Msun$, $T = 60\,\mathrm{K}$}
\begin{align}
    \rho_c &= \left(\frac{c_s^3}{G^{3/2} M}\right)^2 \approx 5 \times 10^{-21}\,\mathrm{g\,cm^{-3}} \\
    n_c &\approx 1300\cmcubed \\
    R &\approx 0.8\text{--}1\pc \quad \textcolor{red}{\neq 0.14\text{--}0.2\pc}
\end{align}

\textbf{Case C: Fix $M = 40\,\Msun$, $R = 0.2\pc$} (required for hydrostatic equilibrium)
\begin{align}
    \rho_c &\sim 10^{-19}\,\mathrm{g\,cm^{-3}}, \quad n_c \sim 3 \times 10^4\cmcubed \\
    c_s &\approx 1.0\kms \\
    T &\approx 280\,\mathrm{K} \quad \textcolor{red}{\neq 60\,\mathrm{K}}
\end{align}

\subsection{Comprehensive Matrix of Initial Conditions}

\begin{table}[h]
\centering
\caption{Matrix of Possible Initial Cloud Configurations}
\small
\begin{tabular}{lcccccl}
\toprule
\textbf{Scenario} & \textbf{T [K]} & \textbf{n$_c$ [cm$^{-3}$]} & \textbf{R [pc]} & \textbf{M [M$_\odot$]} & \textbf{Equilibrium?} & \textbf{Notes} \\
\midrule
\multicolumn{7}{l}{\textit{Literature approaches}} \\
Ballone uniform & 60 & $10^4$ (avg) & 0.14 & 40 & No & Matches obs. \\
Oka Gaussian & -- & variable & $\sigma=0.2$ & 1000$\times$1M$_\odot$ & Virialized & N-body \\
\midrule
\multicolumn{7}{l}{\textit{Bonnor-Ebert at T=60K}} \\
BE (match R) & 60 & $8.7\times10^4$ & 0.14 & 2--3 & \checkmark hydro & M too small \\
BE (match M) & 60 & 1300 & 0.8 & 40 & \checkmark hydro & R too large \\
\midrule
\multicolumn{7}{l}{\textit{Bonnor-Ebert matching M and R}} \\
BE (match M,R) & 280 & $3\times10^4$ & 0.2 & 40 & \checkmark hydro & T too high \\
\midrule
\multicolumn{7}{l}{\textit{KI2000 thermal equilibrium}} \\
KI2000 CNM & 15 & $10^4$ & 0.04 & 0.5 & \checkmark thermal & M,R too small \\
KI2000 unstable & 60 & 1 & 3 & 100 & Unstable & Too diffuse \\
\midrule
\multicolumn{7}{l}{\textit{Turbulent support}} \\
Turbulent cloud & 60 & $10^4$ & 0.2 & 40 & Virial & $\sigma_v \sim 5\kms$ \\
\bottomrule
\end{tabular}
\end{table}

\subsection{Fundamental Conclusion}

\textbf{The observed cloud (M = 40 M$_\odot$, R $\sim$ 0.2 pc, T = 60 K) is NOT in equilibrium.}

The virial analysis confirms this:
\begin{equation}
    \frac{M_\mathrm{virial}}{M_\mathrm{actual}} = \frac{\sigma_v^2 R}{G M} \sim 100
\end{equation}

This indicates the cloud is dominated by:
\begin{itemize}
    \item Tidal forces from the IMBH
    \item Bulk kinetic energy from the encounter
    \item NOT thermal or turbulent pressure support
\end{itemize}

\subsection{Recommended Approaches by Simulation Goal}

\begin{table}[h]
\centering
\caption{Best Initial Conditions by Scientific Goal}
\begin{tabular}{lll}
\toprule
\textbf{Goal} & \textbf{Recommended Setup} & \textbf{Rationale} \\
\midrule
\multirow{2}{*}{Match observations} & Uniform sphere & Simple, directly matches \\
 & $M=40\Msun$, $R=0.14\pc$, $T=60\mathrm{K}$ & observed parameters \\
\midrule
\multirow{2}{*}{Study tidal dynamics} & Gaussian (Oka) & Smooth density gradient, \\
 & $\sigma_r = 0.2\pc$, virialized & shows tidal stretching \\
\midrule
\multirow{3}{*}{Track thermal evolution} & KI2000 equilibrium & Self-consistent cooling, \\
 & $T_\mathrm{init} = 15\mathrm{K}$, $n = 10^4\cmcubed$ & watch shock heating \\
 & (smaller $M$, $R$) & to observed 60 K \\
\midrule
\multirow{2}{*}{Hydrostatic BE sphere} & BE at $T=60\mathrm{K}$ & Analytically tractable, \\
 & Accept $M \sim 2\Msun$ or $R \sim 0.8\pc$ & but doesn't match obs. \\
\midrule
\multirow{2}{*}{Realistic turbulent} & Turbulent cloud & Turbulent pressure \\
 & $\sigma_v = 5\kms$, $M=40\Msun$ & supports larger mass \\
\bottomrule
\end{tabular}
\end{table}

\subsection{If Using KI2000 Equilibrium: Best-Fit Parameters}

For a KI2000-equilibrium cloud that will be heated to $\sim 60\,\mathrm{K}$ during compression:

\begin{table}[h]
\centering
\caption{KI2000 Equilibrium Initial Conditions (Pre-Compression)}
\begin{tabular}{lll}
\toprule
\textbf{Parameter} & \textbf{Value} & \textbf{Constraint} \\
\midrule
Temperature & $T_\mathrm{init} = 15\,\mathrm{K}$ & KI2000 equilibrium at high $n$ \\
Central density & $n_c = 10^4$--$10^5\cmcubed$ & KI2000 stable CNM \\
Sound speed & $c_s = 0.23\kms$ & At 15 K \\
BE radius & $R = 0.02$--$0.05\pc$ & For $\xi_s = 5$ \\
BE mass & $M = 0.1$--$1\,\Msun$ & Much smaller than 40 M$_\odot$ \\
\bottomrule
\end{tabular}
\end{table}

\textbf{Implication}: A KI2000-equilibrium cloud is much more compact than the observed cloud.
You would need $\sim 40$--$400$ such ``cores'' to make up the total 40 M$_\odot$.

Alternatively, embed a small equilibrium core in a larger, non-equilibrium envelope.

%===============================================================================
\section{Shock Heating During Tidal Compression}
\label{sec:shock_heating}
%===============================================================================

This section derives from first principles how tidal compression heats a
cold ($T \sim 15\,\mathrm{K}$) KI2000-equilibrium cloud to the observed
$T = 60\,\mathrm{K}$.

\subsection{The Heating Problem}

\textbf{Key question}: If the initial cloud is in KI2000 equilibrium at
$T_i \approx 15\,\mathrm{K}$, what physical process heats it to $T_f = 60\,\mathrm{K}$?

The answer involves three stages:
\begin{enumerate}
    \item \textbf{Tidal compression} creates converging flows
    \item \textbf{Shocks} convert kinetic energy to thermal energy
    \item \textbf{Turbulent dissipation} maintains elevated temperature against cooling
\end{enumerate}

\subsection{Stage 1: Adiabatic Compression Estimate}

For purely adiabatic compression, the temperature scales as:
\begin{equation}
    \frac{T_2}{T_1} = \left(\frac{\rho_2}{\rho_1}\right)^{\gamma - 1}
\end{equation}

\subsubsection{Required Compression Ratio}
To heat from $T_1 = 15\,\mathrm{K}$ to $T_2 = 60\,\mathrm{K}$ (factor of 4):

\begin{table}[h]
\centering
\caption{Compression Ratio Required for $T: 15 \to 60\,\mathrm{K}$}
\begin{tabular}{lcc}
\toprule
\textbf{Gas Type} & $\gamma$ & \textbf{Compression Ratio} $\rho_2/\rho_1$ \\
\midrule
Monatomic (ideal) & 5/3 & $4^{3/2} = 8$ \\
Diatomic (H$_2$) & 7/5 & $4^{5/2} = 32$ \\
Isothermal limit & 1 & $\infty$ \\
\bottomrule
\end{tabular}
\end{table}

\subsubsection{Observed Compression}
From observations:
\begin{align}
    n_\mathrm{initial} &\sim 10^4\cmcubed \quad \text{(pre-encounter estimate)} \\
    n_\mathrm{final} &\sim 10^{6.5}\cmcubed \quad \text{(observed dense clump)} \\
    \text{Compression ratio} &= 10^{2.5} \approx 300
\end{align}

For this compression ratio, adiabatic heating would give:
\begin{align}
    \gamma = 5/3: \quad T_f &= 15\,\mathrm{K} \times 300^{2/3} = 15 \times 45 \approx 675\,\mathrm{K} \\
    \gamma = 7/5: \quad T_f &= 15\,\mathrm{K} \times 300^{2/5} = 15 \times 9.5 \approx 143\,\mathrm{K}
\end{align}

\textbf{Problem}: Adiabatic compression predicts $T \sim 150$--$700\,\mathrm{K}$,
much hotter than observed 60~K. This means \textbf{cooling is significant} during compression.

\subsection{Stage 2: Shock Heating and Cooling}

\subsubsection{Tidal Compression Velocity}
The tidal acceleration at distance $r$ from the IMBH is:
\begin{equation}
    a_\mathrm{tidal} \sim \frac{G M_\mathrm{BH} R_\mathrm{cloud}}{r^3}
\end{equation}

For $M_\mathrm{BH} = 10^5\,\Msun$, $R_\mathrm{cloud} = 0.2\pc$, $r_\mathrm{peri} = 0.5\pc$:
\begin{equation}
    a_\mathrm{tidal} \sim \frac{4.3 \times 10^{-3} \times 10^5 \times 0.2}{0.5^3}
    \approx 690\,(\kms)^2\,\mathrm{pc}^{-1}
\end{equation}

The compression timescale:
\begin{equation}
    t_\mathrm{compress} \sim \sqrt{\frac{R_\mathrm{cloud}}{a_\mathrm{tidal}}}
    \approx \sqrt{\frac{0.2}{690}}\,\mathrm{pc}\,(\kms)^{-1} \approx 17{,}000\,\mathrm{yr}
\end{equation}

The compression velocity:
\begin{equation}
    v_\mathrm{compress} \sim \frac{R_\mathrm{cloud}}{t_\mathrm{compress}}
    \approx \frac{0.2\pc}{17{,}000\,\mathrm{yr}} \approx 11\kms
\end{equation}

\subsubsection{Mach Number of Compression}
For pre-shock gas at $T = 15\,\mathrm{K}$:
\begin{equation}
    c_s = \sqrt{\frac{k_B T}{\mu m_H}} = 0.23\kms
\end{equation}

Mach number:
\begin{equation}
    \mathcal{M} = \frac{v_\mathrm{compress}}{c_s} = \frac{11}{0.23} \approx 48
\end{equation}

This is a \textbf{highly supersonic} compression, indicating strong shocks.

\subsubsection{Post-Shock Temperature (Rankine-Hugoniot)}
For a strong adiabatic shock ($\mathcal{M} \gg 1$):
\begin{equation}
    \frac{T_\mathrm{post}}{T_\mathrm{pre}} \approx \frac{2\gamma(\gamma-1)}{(\gamma+1)^2} \mathcal{M}^2
\end{equation}

For $\gamma = 5/3$ and $\mathcal{M} = 48$:
\begin{equation}
    T_\mathrm{post} = 15\,\mathrm{K} \times 0.31 \times 48^2 \approx 10{,}700\,\mathrm{K}
\end{equation}

The immediate post-shock temperature is $\sim 10^4\,\mathrm{K}$!

\subsubsection{Rapid Cooling After Shock}
At high temperature and density, cooling is extremely efficient.
The cooling timescale:
\begin{equation}
    t_\mathrm{cool} \sim \frac{k_B T}{n \Lambda(T) (\gamma - 1)}
\end{equation}

For $n = 10^5\cmcubed$, $T = 10^4\,\mathrm{K}$, $\Lambda \sim 10^{-22}\,\mathrm{erg\,cm^3\,s^{-1}}$:
\begin{equation}
    t_\mathrm{cool} \sim \frac{1.38 \times 10^{-16} \times 10^4}{10^5 \times 10^{-22} \times 0.67}
    \approx 2 \times 10^5\,\mathrm{s} \approx 6\,\mathrm{yr}
\end{equation}

\textbf{Key insight}: The cooling time ($\sim 6$ yr) is \textbf{much shorter} than
the encounter timescale ($\sim 10^5$ yr). The gas cools rapidly after each shock.

\subsection{Stage 3: Turbulent Heating Equilibrium}

The observed 60~K is NOT the immediate post-shock temperature, but rather
an \textbf{equilibrium} between continuous turbulent heating and radiative cooling.

\subsubsection{Origin of Turbulence}
The tidal encounter imparts bulk kinetic energy to the cloud.
From virial analysis:
\begin{equation}
    \sigma_v \approx 22\kms \quad \text{(observed velocity dispersion)}
\end{equation}

This supersonic turbulence drives shocks throughout the cloud.

\subsubsection{Turbulent Heating Rate}
The volumetric heating rate from turbulent dissipation:
\begin{equation}
    \Gamma_\mathrm{turb} \sim \frac{\rho \sigma_v^3}{L}
\end{equation}
where $L$ is the driving scale (cloud size).

For $\rho = n \mu m_H = 10^5 \times 3.9 \times 10^{-24} = 3.9 \times 10^{-19}\,\mathrm{g\,cm^{-3}}$,
$\sigma_v = 22\kms = 2.2 \times 10^6\,\mathrm{cm\,s^{-1}}$,
$L = 0.1\pc = 3.1 \times 10^{17}\,\mathrm{cm}$:
\begin{equation}
    \Gamma_\mathrm{turb} = \frac{3.9 \times 10^{-19} \times (2.2 \times 10^6)^3}{3.1 \times 10^{17}}
    \approx 1.3 \times 10^{-17}\,\mathrm{erg\,cm^{-3}\,s^{-1}}
\end{equation}

\subsubsection{Cooling Rate at Equilibrium}
The radiative cooling rate:
\begin{equation}
    \Lambda_\mathrm{cool} = n^2 \Lambda(T)
\end{equation}

At equilibrium, $\Gamma_\mathrm{turb} = \Lambda_\mathrm{cool}$:
\begin{equation}
    \Lambda(T_\mathrm{eq}) = \frac{\Gamma_\mathrm{turb}}{n^2}
    = \frac{1.3 \times 10^{-17}}{(10^5)^2} = 1.3 \times 10^{-27}\,\mathrm{erg\,cm^3\,s^{-1}}
\end{equation}

\subsubsection{Equilibrium Temperature}
From the KI2000 cooling function, $\Lambda(T) \sim 10^{-27}\,\mathrm{erg\,cm^3\,s^{-1}}$
corresponds to:
\begin{equation}
    \boxed{T_\mathrm{eq} \approx 50\text{--}100\,\mathrm{K}}
\end{equation}

\textbf{This matches the observed 60~K!}

\subsection{Summary: The Complete Heating Picture}

\begin{table}[h]
\centering
\caption{Evolution of Cloud Temperature During IMBH Encounter}
\begin{tabular}{lccl}
\toprule
\textbf{Stage} & \textbf{Timescale} & \textbf{Temperature} & \textbf{Physics} \\
\midrule
1. Initial (KI2000 eq.) & -- & 15 K & Thermal equilibrium \\
2. Shock heating & $\sim 10^4$ yr & $\to 10^4$ K & R-H jump conditions \\
3. Rapid cooling & $\sim 10$ yr & $\to 15$ K & Radiative cooling \\
4. Turbulent equilibrium & $\sim 10^5$ yr & $\sim 60$ K & $\Gamma_\mathrm{turb} = \Lambda_\mathrm{cool}$ \\
\bottomrule
\end{tabular}
\end{table}

\begin{figure}[h]
\centering
\fbox{\parbox{0.9\textwidth}{
\textbf{Physical Picture:}\\[0.5em]
The tidal encounter converts gravitational potential energy into bulk kinetic energy
($\sigma_v \sim 22\kms$). This supersonic turbulence continuously dissipates into
heat via shocks. The balance between turbulent heating and radiative cooling
maintains $T \approx 60\,\mathrm{K}$, elevated above the KI2000 equilibrium
($\sim 10\,\mathrm{K}$) but far below the instantaneous post-shock temperature
($\sim 10^4\,\mathrm{K}$).
}}
\end{figure}

\subsection{Implications for SPH Simulation Setup}

Based on this analysis, the recommended simulation approach is:

\begin{enumerate}
    \item \textbf{Initial condition}: KI2000-equilibrium cloud at $T_i \approx 15\,\mathrm{K}$
    \item \textbf{Cooling}: Include KI2000 cooling function
    \item \textbf{Resolution}: Sufficient to resolve shocks ($h \ll R_\mathrm{cloud}$)
    \item \textbf{Expected outcome}: Cloud naturally evolves to $T \sim 60\,\mathrm{K}$
          via turbulent heating equilibrium
\end{enumerate}

\begin{table}[h]
\centering
\caption{Self-Consistent Initial Condition for Shock Heating Simulation}
\begin{tabular}{lll}
\toprule
\textbf{Parameter} & \textbf{Value} & \textbf{Rationale} \\
\midrule
\multicolumn{3}{l}{\textit{Initial cloud (pre-encounter)}} \\
$T_\mathrm{init}$ & $15\,\mathrm{K}$ & KI2000 equilibrium at high $n$ \\
$n_\mathrm{init}$ & $10^4\cmcubed$ & Dense CNM \\
$M_\mathrm{cloud}$ & $40\,\Msun$ & Observed (but distributed) \\
$R_\mathrm{cloud}$ & $0.5$--$1\pc$ & Larger pre-compression \\
\midrule
\multicolumn{3}{l}{\textit{Physics modules}} \\
EOS & Ideal gas, $\gamma = 5/3$ & For shock heating \\
Cooling & KI2000 & Thermal equilibrium \\
Self-gravity & ON (weak) & Secondary to tidal forces \\
\midrule
\multicolumn{3}{l}{\textit{Expected post-encounter state}} \\
$T_\mathrm{final}$ & $\sim 60\,\mathrm{K}$ & Turbulent equilibrium \\
$n_\mathrm{final}$ & $\sim 10^{6.5}\cmcubed$ & Tidal compression \\
$\sigma_v$ & $\sim 20\kms$ & From tidal stretching \\
\bottomrule
\end{tabular}
\end{table}

\textbf{Key prediction}: If you start with a cold ($T = 15\,\mathrm{K}$) cloud
and include proper cooling physics, the tidal encounter will naturally produce
a cloud at $T \approx 60\,\mathrm{K}$ through the turbulent heating mechanism.
This is a self-consistent, predictive simulation setup.

\subsection{Deriving Pre-Encounter Cloud Parameters}

Working backwards from the observed post-encounter state to derive the
pre-encounter initial conditions.

\subsubsection{The Mass-Temperature Constraint at T = 15 K}

For a Bonnor-Ebert sphere at $T = 15\,\mathrm{K}$:
\begin{equation}
    c_s(15\,\mathrm{K}) = \sqrt{\frac{k_B \times 15}{\mu m_H}} = 0.073\kms
\end{equation}

The Jeans mass scales as $M_J \propto c_s^3 / \sqrt{\rho}$.
Compared to $T = 60\,\mathrm{K}$:
\begin{equation}
    \frac{M_J(15\,\mathrm{K})}{M_J(60\,\mathrm{K})} = \left(\frac{c_s(15)}{c_s(60)}\right)^3
    = \left(\frac{0.073}{0.46}\right)^3 = 0.004
\end{equation}

\textbf{Critical finding}: At fixed density, a cold cloud (15~K) has only
\textbf{0.4\%} of the thermal Jeans mass of a warm cloud (60~K)!

For $n = 10^4\cmcubed$ at $T = 15\,\mathrm{K}$:
\begin{equation}
    M_\mathrm{BE}(T=15\,\mathrm{K}, n=10^4) \approx 0.02\,\Msun
\end{equation}

\textbf{Problem}: To get $M = 40\,\Msun$ in thermal equilibrium at 15~K would require
$\sim 2000$ Bonnor-Ebert cores!

\subsubsection{Resolution: Turbulent Support}

Molecular clouds are NOT supported by thermal pressure alone. They are supported
by \textbf{supersonic turbulence}.

For virial equilibrium with turbulent support:
\begin{equation}
    \sigma_v^2 \sim \frac{G M}{R}
\end{equation}

For $M = 40\,\Msun$, $R = 0.5\pc$:
\begin{equation}
    \sigma_v \sim \sqrt{\frac{4.3 \times 10^{-3} \times 40}{0.5}} \approx 0.6\kms
\end{equation}

The turbulent Mach number:
\begin{equation}
    \mathcal{M}_\mathrm{turb} = \frac{\sigma_v}{c_s(15\,\mathrm{K})} = \frac{0.6}{0.073} \approx 8
\end{equation}

This is \textbf{highly supersonic}, consistent with observed molecular cloud turbulence.

\subsubsection{Pre-Encounter Cloud Parameters}

From compression ratio and shock heating analysis:

\begin{align}
    \text{Compression ratio:} \quad \frac{n_\mathrm{final}}{n_\mathrm{initial}} &\sim 300 \\
    n_\mathrm{initial} &\sim \frac{10^{6.5}}{300} \sim 10^4\cmcubed
\end{align}

For a uniform sphere with $M = 40\,\Msun$ and $n = 10^3$--$10^4\cmcubed$:
\begin{align}
    R_\mathrm{initial} &= \left(\frac{3M}{4\pi n \mu m_H}\right)^{1/3} \\
    &\approx 0.3\text{--}0.7\pc
\end{align}

\subsection{Updated Matrix: Two-Phase Approach}

The preferred approach uses gravitational collapse to develop pre-encounter
turbulence naturally, then the IMBH encounter to reach the observed state.

\begin{table}[h]
\centering
\caption{\textcolor{blue}{Preferred Initial Conditions: Two-Phase Approach}}
\small
\begin{tabular}{lll}
\toprule
\textbf{Parameter} & \textbf{Value} & \textbf{Constraint} \\
\midrule
\multicolumn{3}{l}{\textit{Phase 1 Initial: Super-critical BE sphere}} \\
$M_\mathrm{cloud}$ & $40\,\Msun$ & Observed total mass \\
$R_\mathrm{init}$ & $\sim 3\pc$ & From BE solution at $T = 15\,\mathrm{K}$ \\
$n_\mathrm{center}$ & $100$--$300\cmcubed$ & BE sphere centrally concentrated \\
$T_\mathrm{init}$ & $15\,\mathrm{K}$ & KI2000 CNM equilibrium \\
$\sigma_{v,\mathrm{init}}$ & \textbf{0} & \textbf{No injected turbulence} \\
Density profile & \textbf{Bonnor-Ebert} & Physically motivated, analytically tractable \\
\midrule
\multicolumn{3}{l}{\textit{After Phase 1 collapse ($\sim 2$--$3$ Myr)}} \\
$R$ & $0.3$--$0.5\pc$ & Collapsed ($\times 10$ compression) \\
$n$ & $\sim 10^4\cmcubed$ & Compressed ($\times 100$ from center) \\
$\sigma_v$ & $1$--$2\kms$ & \textbf{From collapse, not injection} \\
\midrule
\multicolumn{3}{l}{\textit{After Phase 2 IMBH encounter}} \\
$\sigma_v$ & $\sim 22\kms$ & From tidal forces \\
$T$ & $\sim 60\,\mathrm{K}$ & Shock heating equilibrium \\
\bottomrule
\end{tabular}
\end{table}

\textbf{Key insight}: Start with a \textbf{larger, more diffuse} cloud. Let
gravitational collapse generate $\sigma_v \sim 1$--$2\kms$ naturally. Then
the IMBH encounter adds $\Delta v \sim 20$--$100\kms$ from tidal forces.

\subsection{Comprehensive Comparison: All Initial Condition Options}

\begin{table}[h]
\centering
\caption{Complete Matrix of Initial Conditions with Two-Phase Assessment}
\footnotesize
\begin{tabular}{lccccccl}
\toprule
\textbf{Scenario} & \textbf{T [K]} & \textbf{n$_c$ [cm$^{-3}$]} & \textbf{R [pc]} & \textbf{M [M$_\odot$]} & \textbf{$\sigma_v$} & \textbf{Turbulence} & \textbf{Verdict} \\
\midrule
\multicolumn{8}{l}{\textit{\textcolor{blue}{Recommended: Super-critical BE sphere + IMBH tidal}}} \\
\textbf{BE sphere $\to$ IMBH} & \textbf{15} & \textbf{100--300} & \textbf{$\sim$3} & \textbf{40} & \textbf{0} & \textbf{IMBH tidal} & \textcolor{blue}{\textbf{Best}} \\
\midrule
\multicolumn{8}{l}{\textit{Alternative: Direct encounter (no pre-collapse)}} \\
Uniform cold & 15 & $10^4$ & 0.3--0.5 & 40 & 0 & From tidal & Simpler \\
Ballone uniform & 60 & $10^4$ & 0.14 & 40 & 0 & From tidal & Matches final \\
\midrule
\multicolumn{8}{l}{\textit{Literature}} \\
Oka Gaussian & -- & variable & 0.2 & 1000 & virial & Pre-set & N-body only \\
\midrule
\multicolumn{8}{l}{\textit{Not recommended}} \\
Injected turbulence & 15 & $10^4$ & 0.5 & 40 & artificial & Artificial & Unphysical \\
BE T=15K, small & 15 & $10^4$ & 0.04 & 0.02 & 0 & -- & M too small \\
\bottomrule
\end{tabular}
\end{table}

\subsection{Physical Justification: Why No Turbulence Injection}

The uniform cold sphere (no injected turbulence) is preferred because:

\begin{enumerate}
    \item \textbf{Observed turbulence is post-encounter}: The $\sigma_v \sim 22\kms$
          velocity dispersion is \textbf{caused by} the tidal encounter, not pre-existing

    \item \textbf{Timescale argument}: The encounter timescale ($t_\mathrm{enc} \sim 0.1\,\mathrm{Myr}$)
          is shorter than the free-fall time ($t_\mathrm{ff} \sim 1\,\mathrm{Myr}$),
          so the cloud does not need to be in equilibrium

    \item \textbf{Tidal dominance}: The tidal radius $r_t \ll R_\mathrm{cloud}$,
          meaning IMBH gravity completely dominates over cloud self-gravity

    \item \textbf{Physically motivated}: Letting tidal forces generate the velocity
          dispersion is more realistic than artificial injection

    \item \textbf{Predictive}: Shock heating from tidal compression will naturally
          produce $T \to 60\,\mathrm{K}$ via turbulent dissipation equilibrium
\end{enumerate}

\subsection{Timescale Analysis: Why Equilibrium is Not Required}

\begin{table}[h]
\centering
\caption{Relevant Timescales for IMBH-Cloud Encounter}
\begin{tabular}{lcl}
\toprule
\textbf{Timescale} & \textbf{Value} & \textbf{Implication} \\
\midrule
Free-fall time & $t_\mathrm{ff} \sim 1\,\mathrm{Myr}$ & Cloud collapse time \\
Encounter time & $t_\mathrm{enc} \sim 0.1\,\mathrm{Myr}$ & Pericenter passage \\
Cooling time & $t_\mathrm{cool} \sim 10\,\mathrm{yr}$ & Rapid thermal equilibration \\
\midrule
\multicolumn{3}{l}{\textbf{Key relation}: $t_\mathrm{enc} < t_\mathrm{ff}$ $\Rightarrow$ encounter happens before collapse} \\
\bottomrule
\end{tabular}
\end{table}

Since $t_\mathrm{enc} \ll t_\mathrm{ff}$, the cloud does not have time to collapse
before the IMBH encounter. Therefore, the initial condition does \textbf{not} need
to be in hydrostatic or virial equilibrium.

\subsection{Collapse-Driven Turbulence: First Principles}

A super-Jeans cloud undergoing gravitational collapse naturally generates
turbulent velocities without any artificial injection.

\subsubsection{How Super-Jeans is the Cloud?}

For $T = 15\,\mathrm{K}$ and $n = 10^4\cmcubed$:
\begin{align}
    M_\mathrm{Jeans} &\approx 0.02\,\Msun \\
    M_\mathrm{cloud} &= 40\,\Msun \\
    \frac{M}{M_J} &\sim 2000 \quad \text{(highly super-Jeans)}
\end{align}

This cloud is \textbf{extremely} gravitationally unstable and will collapse.

\subsubsection{Free-Fall Velocity}

The characteristic velocity generated during collapse:
\begin{equation}
    v_\mathrm{ff} = \sqrt{\frac{2 G M}{R}}
\end{equation}

\begin{table}[h]
\centering
\caption{Free-Fall Velocity vs Cloud Radius ($M = 40\,\Msun$)}
\begin{tabular}{cccc}
\toprule
\textbf{R [pc]} & \textbf{$v_\mathrm{ff}$ [km/s]} & \textbf{Mach (T=15K)} & \textbf{State} \\
\midrule
1.0 & 0.6 & 8 & Initial (diffuse) \\
0.5 & 0.8 & 11 & Collapsing \\
0.3 & 1.1 & 15 & Collapsing \\
0.2 & 1.3 & 18 & Collapsed \\
0.1 & 1.8 & 25 & Dense core \\
\bottomrule
\end{tabular}
\end{table}

\textbf{Key insight}: Collapse from $R = 1\pc$ to $R \sim 0.2\pc$ naturally
generates $\sigma_v \sim 1$--$2\kms$!

\subsubsection{Mechanisms for Turbulence Generation}

During gravitational collapse, several mechanisms create turbulent motions:
\begin{enumerate}
    \item \textbf{Converging infall}: Material falling from different directions
          creates shocks at the center
    \item \textbf{Gravitational fragmentation}: The cloud fragments into clumps,
          each with its own velocity
    \item \textbf{Angular momentum}: Any initial rotation gets amplified
          (conservation of angular momentum)
    \item \textbf{Thermal instability}: KI2000 cooling creates density contrasts
          that seed fragmentation
\end{enumerate}

\subsubsection{Free-Fall Timescale}

\begin{equation}
    t_\mathrm{ff} = \sqrt{\frac{3\pi}{32 G \rho}}
\end{equation}

\begin{table}[h]
\centering
\caption{Free-Fall Time vs Density}
\begin{tabular}{cc}
\toprule
\textbf{n [cm$^{-3}$]} & \textbf{$t_\mathrm{ff}$ [Myr]} \\
\midrule
100 & 3.3 \\
300 & 1.9 \\
1000 & 1.0 \\
3000 & 0.6 \\
10000 & 0.3 \\
\bottomrule
\end{tabular}
\end{table}

\subsection{Two-Phase Simulation Strategy}

The simulation proceeds in two phases: (1) collapse to develop turbulence,
then (2) IMBH encounter.

\subsubsection{Phase 1: Super-Critical Bonnor-Ebert Sphere}

\textbf{Why Bonnor-Ebert instead of Uniform Sphere?}

For a \textbf{critical} Bonnor-Ebert sphere at $T = 15\,\mathrm{K}$:
\begin{equation}
    M_\mathrm{crit} \approx 1.18 \frac{c_s^4}{G^{3/2} \sqrt{P_\mathrm{ext}}} \sim 1\text{--}2\,\Msun
\end{equation}

A $40\,\Msun$ BE sphere at $T = 15\,\mathrm{K}$ is therefore \textbf{highly super-critical}:
\begin{equation}
    \frac{M}{M_\mathrm{crit}} \sim 20\text{--}40
\end{equation}

This means:
\begin{itemize}
    \item The BE sphere is \textbf{gravitationally unstable} and will collapse
    \item The BE profile provides a \textbf{physically motivated} initial density distribution
    \item No artificial perturbations needed --- the instability is built-in
\end{itemize}

\textbf{Initial condition} (super-critical BE sphere):

\begin{table}[h]
\centering
\caption{\textcolor{blue}{Recommended Initial Condition: Super-Critical Bonnor-Ebert Sphere}}
\begin{tabular}{lll}
\toprule
\textbf{Parameter} & \textbf{Value} & \textbf{Rationale} \\
\midrule
$M_\mathrm{cloud}$ & $40\,\Msun$ & Observed mass \\
$R_\mathrm{cloud}$ & $\sim 3\pc$ & From BE solution at $T = 15\,\mathrm{K}$ \\
$n_\mathrm{center}$ & $100$--$300\cmcubed$ & Centrally concentrated \\
$T$ & $15\,\mathrm{K}$ & KI2000 CNM equilibrium \\
$\sigma_{v,\mathrm{init}}$ & \textbf{0} & No injected turbulence \\
Density profile & \textbf{Bonnor-Ebert} & Analytically tractable, physical \\
\bottomrule
\end{tabular}
\end{table}

\textbf{Advantages of BE sphere over uniform sphere:}
\begin{itemize}
    \item \textbf{Centrally concentrated}: More realistic density profile
    \item \textbf{Analytically defined}: Lane-Emden solution, reproducible
    \item \textbf{Initially in hydrostatic equilibrium}: Clean starting point
    \item \textbf{Larger initial radius}: More time for IMBH tidal forces to act
    \item \textbf{Compression ratio}: $R: 3\pc \to 0.3\pc$ gives density increase $\times 1000$
          (matches observed $n \sim 10^{6.5}\cmcubed$!)
\end{itemize}

\textbf{Phase 1 evolution} (with IMBH present from start):
\begin{enumerate}
    \item $t = 0$: Super-critical BE sphere ($R \sim 3\pc$, $n_c \sim 200\cmcubed$, $\sigma_v = 0$)
    \item $t \sim 0.5\,t_\mathrm{ff}$: Cloud contracting under self-gravity + tidal forces
          shaping the collapse, infall velocities developing
    \item $t \sim 1\,t_\mathrm{ff}$: Cloud collapsed to $R \sim 0.3$--$0.5\pc$,
          $n \sim 10^4\cmcubed$, $\sigma_v \sim 1$--$2\kms$
\end{enumerate}

\textbf{Phase 1 duration}: $\sim 1$--$2\,\mathrm{Myr}$ (one free-fall time at initial density)

\subsubsection{Key Insight: IMBH as the Perturbation Source}

The traditional approach to simulating cloud collapse requires adding artificial
density perturbations (e.g., 5--10\% random fluctuations) to seed fragmentation.
However, in our IMBH-cloud interaction scenario:

\begin{equation}
    \boxed{\text{IMBH tidal forces} \Rightarrow \text{Density perturbations (no artificial seeding)}}
\end{equation}

\textbf{Physical mechanism}:
\begin{enumerate}
    \item The tidal acceleration varies across the cloud:
          $\Delta a_\mathrm{tidal} \sim (G M_\mathrm{BH} / r^3) \times R_\mathrm{cloud}$
    \item Material closer to IMBH experiences stronger acceleration
    \item This creates a velocity gradient $\to$ density gradient
    \item The cloud is stretched along the IMBH direction and compressed perpendicular to it
    \item These tidal-induced density gradients seed gravitational fragmentation
\end{enumerate}

\textbf{Order of magnitude}:
For $M_\mathrm{BH} = 10^5\,\Msun$ at $r = 10\pc$ from a cloud with $R = 3\pc$:
\begin{align}
    \Delta a_\mathrm{tidal} &\sim \frac{G M_\mathrm{BH}}{r^3} R_\mathrm{cloud}
    \sim \frac{4.3 \times 10^{-3} \times 10^5}{1000} \times 3 \approx 1.3\,(\kms)^2\,\mathrm{pc}^{-1}
\end{align}

Over one free-fall time ($\sim 3\,\mathrm{Myr}$ at $n \sim 200\cmcubed$), this creates velocity differences of:
\begin{equation}
    \Delta v \sim \Delta a \times t_\mathrm{ff} \times R \sim 1.3 \times 3 \times 3 \approx 12\kms
\end{equation}

This is \textbf{much larger} than the sound speed ($c_s \sim 0.23\kms$ at 15~K),
confirming that tidal perturbations are highly supersonic and will dominate over
any thermal or turbulent effects.

\subsubsection{Simulation Strategy: IMBH Present from Start}

\textbf{Important}: Since the IMBH provides the density perturbations, it should be
present from $t = 0$, not introduced later. Options:

\begin{enumerate}
    \item \textbf{IMBH at large initial distance} ($r_0 \sim 5$--$10\pc$): Cloud initially
          feels weak tidal field. As IMBH approaches on flyby trajectory, tidal forces
          strengthen and trigger collapse/fragmentation. \textit{Most physical.}

    \item \textbf{Cloud falling toward IMBH}: Cloud starts at large distance and falls
          radially toward IMBH. Tidal effects grow as $r$ decreases. \textit{Simpler orbit.}

    \item \textbf{IMBH on fixed orbit}: IMBH follows prescribed trajectory (e.g., from
          Oka's N-body fit) while cloud responds. \textit{Matches observation geometry.}
\end{enumerate}

\subsubsection{Phase 2: Close IMBH Encounter}

Once the cloud has been shaped by tidal forces during approach:

\begin{enumerate}
    \item \textbf{Place IMBH}: At distance $\sim 10\pc$, on flyby trajectory
    \item \textbf{Cloud state at encounter start}:
    \begin{itemize}
        \item $R \sim 0.3$--$0.5\pc$ (collapsed)
        \item $n \sim 10^4\cmcubed$
        \item $\sigma_v \sim 1$--$2\kms$ (from collapse)
        \item $T \sim 15\,\mathrm{K}$ (KI2000 equilibrium)
    \end{itemize}
    \item \textbf{Tidal forces add}: $\Delta v \sim 20$--$100\kms$
    \item \textbf{Final state}: $\sigma_v \sim 22\kms$, $T \sim 60\,\mathrm{K}$, $n \sim 10^{6.5}\cmcubed$
\end{enumerate}

\subsection{Complete Evolution Timeline}

\begin{table}[h]
\centering
\caption{Simulation Timeline: IMBH Present from Start (Super-Critical BE Sphere)}
\begin{tabular}{lcccccc}
\toprule
\textbf{Phase} & \textbf{Time} & \textbf{R [pc]} & \textbf{n$_c$ [cm$^{-3}$]} & \textbf{T [K]} & \textbf{$\sigma_v$ [km/s]} & \textbf{IMBH dist [pc]} \\
\midrule
Initial (BE sphere) & 0 & $\sim 3$ & 100--300 & 15 & 0 & 10--15 \\
Tidal shaping & 1 Myr & 1.5--2 & $10^3$ & 15 & 0.5--1 & 5--10 \\
Pre-pericenter & 2--3 Myr & 0.3--0.5 & $\sim 10^4$ & 15 & \textbf{1--2} & 1--2 \\
\midrule
Pericenter & +0.1 Myr & stretching & $10^5$--$10^6$ & $10^4$ (shock) & 50--100 & $<1$ \\
Post-encounter & +0.2 Myr & 0.3 & $10^{6.5}$ & \textbf{60} & \textbf{22} & receding \\
\bottomrule
\end{tabular}
\end{table}

\textbf{Key point}: The IMBH is present throughout, starting at large distance.
As it approaches, tidal forces progressively shape the cloud, naturally creating
the density gradients that seed collapse and fragmentation.

\begin{figure}[h]
\centering
\fbox{\parbox{0.9\textwidth}{
\textbf{Summary}: Start with a \textbf{super-critical Bonnor-Ebert sphere} ($M = 40\,\Msun$,
$R \sim 3\pc$, $T = 15\,\mathrm{K}$, $n_c \sim 200\cmcubed$) with \textbf{zero velocity}.
The BE sphere is highly super-critical ($M/M_\mathrm{crit} \sim 20$--$40$), so it is
gravitationally unstable. No density perturbations needed --- the IMBH tidal forces will
create density gradients naturally. The cloud collapses from $R \sim 3\pc$ to $R \sim 0.3\pc$,
giving a density enhancement of $\times 1000$ (matching observed $n \sim 10^{6.5}\cmcubed$!).
The IMBH encounter adds $\Delta v \sim 20$--$100\kms$ from tidal forces, producing the
observed state with $T \sim 60\,\mathrm{K}$ and $\sigma_v \sim 22\kms$.
}}
\end{figure}

%===============================================================================
\section{Failure of Isothermal Approach and Switch to $\gamma = 5/3$ Polytrope}
\label{sec:polytrope}
%===============================================================================

\subsection{The Isothermal Instability Problem}

Numerical simulations with the isothermal equation of state ($\gamma \to 1$) revealed
a fundamental problem: \textbf{no stable equilibrium exists} for our cloud parameters.

\subsubsection{Bonnor-Ebert Critical Mass}

For an isothermal Bonnor-Ebert sphere, the maximum stable mass is:
\begin{equation}
    M_\mathrm{crit} = 1.18 \frac{c_s^4}{\sqrt{G^3 P_\mathrm{ext}}}
\end{equation}

At $T = 15\,\mathrm{K}$ with $c_s = 0.23\kms$ and $P_\mathrm{ext}/k_B = 75\,\mathrm{K\,cm^{-3}}$:
\begin{equation}
    M_\mathrm{crit} \approx 1\text{--}2\,\Msun
\end{equation}

Our $40\,\Msun$ cloud is therefore \textbf{20--40 times super-critical}.

\subsubsection{Numerical Evidence: Runaway Collapse}

Simulations with $\gamma = 1.0001$ showed:
\begin{itemize}
    \item Initial cloud: $R = 2.4\pc$, $n_c = 200\cmcubed$
    \item After $t \sim 0.18\,\mathrm{Myr}$ (one free-fall time): timestep collapsed to $\Delta t \sim 10^{-10}\,\mathrm{Myr}$
    \item Central density runaway toward singularity
    \item Simulation became computationally intractable
\end{itemize}

This is the \textbf{expected physical behavior} for a super-critical isothermal sphere---it
simply cannot exist in equilibrium.

\subsection{Solution: Polytropic EOS with $\gamma = 5/3$}

We switch to an adiabatic equation of state:
\begin{equation}
    P = K \rho^\gamma \quad \text{with} \quad \gamma = 5/3
\end{equation}

\subsubsection{Stronger Pressure Response}

The key difference between isothermal and adiabatic compression:
\begin{align}
    \text{Isothermal:} \quad P &\propto \rho \quad \text{(weak)} \\
    \text{Adiabatic:} \quad P &\propto \rho^{5/3} \quad \text{(strong)}
\end{align}

When compressed, adiabatic gas builds up pressure faster, providing self-regulating
support against collapse.

\subsubsection{Lane-Emden Polytrope (n = 1.5)}

For $\gamma = 5/3$, the polytropic index is $n = 1/(\gamma - 1) = 1.5$.
The Lane-Emden equation has a well-defined, finite solution with:
\begin{itemize}
    \item Dimensionless radius: $\xi_1 = 3.65375$
    \item Mass concentration: $-\xi^2 (d\theta/d\xi)|_{\xi_1} = 2.71406$
    \item Stable equilibrium for \textbf{any} mass
\end{itemize}

\subsubsection{Adiabatic Heating During Compression: Detailed Derivation}

For an adiabatic process (no heat exchange with surroundings), the first law of
thermodynamics combined with the ideal gas law gives a fundamental relationship
between temperature and density.

\paragraph{Step 1: Adiabatic Condition}
For an adiabatic process, pressure and volume satisfy:
\begin{equation}
    P V^\gamma = \text{constant}
\end{equation}
Since $V \propto 1/\rho$ for a fixed mass, this becomes:
\begin{equation}
    P \rho^{-\gamma} = \text{constant} \quad \Rightarrow \quad P \propto \rho^\gamma
\end{equation}

\paragraph{Step 2: Ideal Gas Law}
For an ideal gas:
\begin{equation}
    P = \frac{\rho k_B T}{\mu m_H}
\end{equation}
where $\mu$ is the mean molecular weight and $m_H$ is the hydrogen mass.

\paragraph{Step 3: Temperature-Density Relation}
Substituting the ideal gas law into the adiabatic condition:
\begin{equation}
    \frac{\rho k_B T}{\mu m_H} \propto \rho^\gamma
\end{equation}
Solving for temperature:
\begin{equation}
    T \propto \rho^{\gamma - 1}
\end{equation}

\paragraph{Step 4: Application to $\gamma = 5/3$}
For a monatomic or adiabatic gas with $\gamma = 5/3$:
\begin{equation}
    \boxed{T \propto \rho^{2/3}}
\end{equation}

If the density increases by a factor $\chi = \rho_2/\rho_1$, the temperature increases by:
\begin{equation}
    \frac{T_2}{T_1} = \chi^{2/3} = \left(\frac{\rho_2}{\rho_1}\right)^{2/3}
\end{equation}

\paragraph{Step 5: Numerical Example}
For a compression factor of $\chi = 10$:
\begin{equation}
    \frac{T_2}{T_1} = 10^{2/3} = 10^{0.667} \approx 4.64
\end{equation}

Starting from $T_1 = 15\,\mathrm{K}$:
\begin{equation}
    T_2 = 15\,\mathrm{K} \times 4.64 \approx 70\,\mathrm{K}
\end{equation}

\begin{table}[h]
\centering
\caption{Adiabatic Heating vs Compression Ratio ($\gamma = 5/3$)}
\begin{tabular}{cccc}
\toprule
\textbf{Compression $\rho_2/\rho_1$} & \textbf{$(\rho_2/\rho_1)^{2/3}$} & \textbf{$T_2/T_1$} & \textbf{$T_2$ (from $T_1 = 15\,\mathrm{K}$)} \\
\midrule
2 & $2^{0.667}$ & 1.59 & 24 K \\
3 & $3^{0.667}$ & 2.08 & 31 K \\
5 & $5^{0.667}$ & 2.92 & 44 K \\
\textbf{10} & $\mathbf{10^{0.667}}$ & \textbf{4.64} & \textbf{70 K} \\
20 & $20^{0.667}$ & 7.37 & 111 K \\
30 & $30^{0.667}$ & 9.65 & 145 K \\
100 & $100^{0.667}$ & 21.5 & 323 K \\
\bottomrule
\end{tabular}
\end{table}

\paragraph{Physical Interpretation}
The $T \propto \rho^{2/3}$ scaling arises because:
\begin{itemize}
    \item When gas is compressed, work is done \textit{on} the gas
    \item This work increases the internal (thermal) energy
    \item For $\gamma = 5/3$, the internal energy $U = \frac{3}{2} N k_B T$
    \item The work-energy balance gives the $\rho^{2/3}$ dependence
\end{itemize}

\textbf{Key insight}: A compression factor of $\sim 10$ produces $T \approx 70\,\mathrm{K}$,
matching the observed $60\,\mathrm{K}$! The observed temperature is a \textbf{natural outcome}
of adiabatic compression, not an input assumption.

\paragraph{Why This Matters for IMBH-Cloud Interaction}
During the tidal encounter:
\begin{enumerate}
    \item Pre-encounter cloud: $R \sim 1\pc$, $n \sim 10^3$--$10^4\cmcubed$, $T \sim 15\,\mathrm{K}$
    \item Tidal compression: Cloud squeezed to $R \sim 0.3\pc$ (volume $\times 0.027$, density $\times 37$)
    \item Post-encounter: Adiabatic heating gives $T \sim 15 \times 37^{0.667} \sim 165\,\mathrm{K}$
    \item With radiative cooling: Temperature settles to observed $\sim 60\,\mathrm{K}$
\end{enumerate}

The $\gamma = 5/3$ model naturally produces temperatures in the right range,
while the isothermal model ($\gamma = 1$) would keep $T = 15\,\mathrm{K}$ regardless
of compression---failing to explain the observed 60 K.

\subsection{Hybrid Approach: $\gamma = 5/3$ with Cooling During Flyby}
\label{sec:hybrid_cooling}

A more physically realistic simulation uses a \textbf{hybrid approach}:
\begin{itemize}
    \item \textbf{Phase 1--2 (Relaxation \& Stability Test)}: $\gamma = 5/3$, cooling OFF
    \item \textbf{Phase 3 (IMBH Flyby)}: $\gamma = 5/3$, cooling ON
\end{itemize}

This captures the essential physics: stable initial conditions with realistic
thermodynamics during the tidal encounter.

\subsubsection{Physics of Cooling-Triggered Isothermal Shocks}

When cooling is enabled during the flyby, a rich shock structure emerges:

\paragraph{Step 1: Tidal Compression}
The IMBH's tidal forces compress the cloud perpendicular to the orbital plane:
\begin{equation}
    a_\mathrm{tidal} \sim \frac{G M_\mathrm{BH} R_\mathrm{cloud}}{r^3}
\end{equation}

\paragraph{Step 2: Adiabatic Heating (Initial Response)}
The compression initially heats the gas adiabatically:
\begin{equation}
    T \propto \rho^{2/3} \quad \Rightarrow \quad T_\mathrm{post-shock} \gg T_\mathrm{pre-shock}
\end{equation}

\paragraph{Step 3: Radiative Cooling}
The heated gas radiates energy via molecular line emission (CO, H$_2$O, etc.):
\begin{equation}
    \dot{E}_\mathrm{cool} = n^2 \Lambda(T)
\end{equation}
where $\Lambda(T)$ is the cooling function (e.g., KI2000).

\paragraph{Step 4: Cooling Time vs Compression Time}
The critical comparison is:
\begin{align}
    t_\mathrm{cool} &\sim \frac{k_B T}{n \Lambda(T) (\gamma - 1)} \\
    t_\mathrm{compress} &\sim \sqrt{\frac{R_\mathrm{cloud}}{a_\mathrm{tidal}}}
\end{align}

\begin{itemize}
    \item If $t_\mathrm{cool} \ll t_\mathrm{compress}$: Gas cools faster than it compresses
          $\Rightarrow$ \textbf{isothermal shock}
    \item If $t_\mathrm{cool} \gg t_\mathrm{compress}$: Gas remains hot
          $\Rightarrow$ \textbf{adiabatic shock}
\end{itemize}

\paragraph{Step 5: Isothermal Shock Formation}
When cooling is efficient, the post-shock gas rapidly loses thermal energy:
\begin{enumerate}
    \item Compression heats gas to $T \sim 10^3$--$10^4\,\mathrm{K}$
    \item Rapid cooling (molecular lines) brings $T$ back down
    \item Loss of thermal pressure support $\Rightarrow$ further collapse
    \item Dense, cold layer forms behind the shock $\Rightarrow$ \textbf{isothermal shock}
\end{enumerate}

\begin{figure}[h]
\centering
\fbox{\parbox{0.9\textwidth}{
\textbf{Isothermal Shock Structure:}
\begin{verbatim}
     Pre-shock       Shock front      Post-shock (cooling zone)     Cooled layer
    ============    =============    ==========================    =============
    T ~ 15 K        T jumps to       T drops from ~10^3 K          T ~ 60 K
    n ~ 10^3        ~10^3-10^4 K     back toward ~60 K             n ~ 10^6
    v ~ 0           (adiabatic)      (radiative cooling)           (compressed)
\end{verbatim}
}}
\end{figure}

\subsubsection{Isothermal Shock Jump Conditions}

For a strong isothermal shock (Mach number $\mathcal{M} \gg 1$):
\begin{align}
    \text{Density jump:} \quad \frac{\rho_2}{\rho_1} &= \mathcal{M}^2 \\
    \text{Velocity jump:} \quad \frac{v_2}{v_1} &= \frac{1}{\mathcal{M}^2} \\
    \text{Temperature:} \quad T_2 &= T_1 \quad \text{(isothermal)}
\end{align}

Compare to adiabatic shock ($\gamma = 5/3$):
\begin{align}
    \frac{\rho_2}{\rho_1} &= \frac{(\gamma+1)\mathcal{M}^2}{(\gamma-1)\mathcal{M}^2 + 2} \to 4 \quad (\mathcal{M} \to \infty)
\end{align}

\textbf{Key difference}: Isothermal shocks can achieve \textbf{arbitrarily high compression}
($\rho_2/\rho_1 = \mathcal{M}^2$), while adiabatic shocks are limited to $\rho_2/\rho_1 \leq 4$.

\subsubsection{Why This Matters for IMBH-Cloud Interaction}

The observed dense clump has $n \sim 10^{6.5}\cmcubed$, requiring compression of
$\sim 100$--$1000\times$ from the initial state. This is \textbf{only possible with
isothermal shocks} (cooling-dominated).

\begin{table}[h]
\centering
\caption{Shock Compression: Adiabatic vs Isothermal}
\begin{tabular}{lcc}
\toprule
\textbf{Property} & \textbf{Adiabatic ($\gamma = 5/3$)} & \textbf{Isothermal (with cooling)} \\
\midrule
Max compression & $\rho_2/\rho_1 \leq 4$ & $\rho_2/\rho_1 = \mathcal{M}^2$ (unlimited) \\
For $\mathcal{M} = 10$ & $\rho_2/\rho_1 = 3.9$ & $\rho_2/\rho_1 = 100$ \\
For $\mathcal{M} = 30$ & $\rho_2/\rho_1 = 4.0$ & $\rho_2/\rho_1 = 900$ \\
Post-shock T & $T_2 \propto \mathcal{M}^2$ (hot) & $T_2 \approx T_1$ (cold) \\
Dense clump? & \textcolor{red}{Cannot form} & \textcolor{blue}{Can form} \\
\bottomrule
\end{tabular}
\end{table}

\subsubsection{Simulation Strategy: Cooling Switch}

\begin{table}[h]
\centering
\caption{\textcolor{blue}{Recommended Hybrid Approach: $\gamma = 5/3$ + Cooling}}
\begin{tabular}{llll}
\toprule
\textbf{Phase} & \textbf{EOS} & \textbf{Cooling} & \textbf{Rationale} \\
\midrule
1. Relaxation & $\gamma = 5/3$ & OFF & Stable equilibrium \\
2. Hydrostatic Test & $\gamma = 5/3$ & OFF & Verify stability \\
3. IMBH Flyby & $\gamma = 5/3$ & \textbf{ON} & Isothermal shocks form \\
\bottomrule
\end{tabular}
\end{table}

\textbf{Expected behavior during Phase 3}:
\begin{enumerate}
    \item Cloud initially in $\gamma = 5/3$ equilibrium ($T \sim 15\,\mathrm{K}$)
    \item IMBH tidal forces compress cloud
    \item Compression heats gas adiabatically ($T \to 100$--$1000\,\mathrm{K}$)
    \item Cooling rapidly removes thermal energy ($t_\mathrm{cool} \sim 10$--$100\,\mathrm{yr}$)
    \item Temperature drops back toward equilibrium ($T \to 60\,\mathrm{K}$)
    \item Loss of pressure support $\Rightarrow$ further collapse
    \item Dense clump forms via isothermal shock ($n \to 10^{6.5}\cmcubed$)
\end{enumerate}

\subsection{Observational Consistency with $\gamma = 5/3$}

The observations constrain the \textbf{post-encounter} state, not the pre-encounter EOS:

\begin{table}[h]
\centering
\caption{Pre- vs Post-Encounter Cloud Properties}
\begin{tabular}{lcc}
\toprule
\textbf{Property} & \textbf{Pre-Encounter} & \textbf{Post-Encounter (Observed)} \\
\midrule
Temperature & $\sim 15\,\mathrm{K}$ & $60\,\mathrm{K}$ \\
Density & $n \sim 10^4\cmcubed$ & $n \sim 10^{6.5}\cmcubed$ \\
Radius & $R \sim 0.5$--$1\pc$ & $R \sim 0.3\pc$ \\
Velocity dispersion & $\sim 1\kms$ & $\sim 22\kms$ \\
Equilibrium & Yes ($\gamma = 5/3$ polytrope) & No (tidally disrupted) \\
\bottomrule
\end{tabular}
\end{table}

The $\gamma = 5/3$ model \textbf{predicts} the observed 60 K temperature through
adiabatic compression, rather than requiring it as an initial condition.

\subsection{Physical Motivation}

Real molecular clouds are not purely isothermal. They have:
\begin{enumerate}
    \item \textbf{Turbulent pressure}: Supersonic turbulence provides effective pressure support,
          equivalent to $\gamma_\mathrm{eff} > 1$
    \item \textbf{Magnetic pressure}: $B$-fields resist compression
    \item \textbf{Finite cooling time}: On short timescales, compression is approximately adiabatic
\end{enumerate}

The $\gamma = 5/3$ EOS effectively captures these non-thermal pressure sources
while keeping the simulation tractable.

\subsection{Comparison: Isothermal vs Adiabatic}

\begin{table}[h]
\centering
\caption{Comparison of Isothermal and Adiabatic EOS for IMBH-Cloud Simulation}
\begin{tabular}{lcc}
\toprule
\textbf{Property} & \textbf{Isothermal ($\gamma = 1$)} & \textbf{Adiabatic ($\gamma = 5/3$)} \\
\midrule
Equilibrium for $M = 40\,\Msun$ & \textcolor{red}{No (super-critical)} & \textcolor{blue}{Yes (Lane-Emden)} \\
Stability & \textcolor{red}{Unstable $\to$ collapse} & \textcolor{blue}{Stable} \\
Pressure response to compression & Weak ($P \propto \rho$) & Strong ($P \propto \rho^{5/3}$) \\
Post-compression temperature & Constant & $T \propto \rho^{2/3}$ \\
Observed 60 K explained by & External heating required & Adiabatic compression \\
Numerical behavior & Runaway collapse & Stable initial condition \\
\bottomrule
\end{tabular}
\end{table}

\subsection{Recommended Polytropic Parameters}

\begin{table}[h]
\centering
\caption{\textcolor{blue}{Recommended Initial Condition: $\gamma = 5/3$ Polytrope}}
\begin{tabular}{lll}
\toprule
\textbf{Parameter} & \textbf{Value} & \textbf{Rationale} \\
\midrule
EOS & $P = K \rho^{5/3}$ & Stable polytrope \\
$M_\mathrm{cloud}$ & $40\,\Msun$ & Observed mass \\
$R_\mathrm{cloud}$ & $0.5$--$1\pc$ & Pre-compression size \\
$n_\mathrm{center}$ & $10^3$--$10^4\cmcubed$ & Lane-Emden profile \\
$T_\mathrm{initial}$ & $15$--$20\,\mathrm{K}$ & Sets polytropic constant $K$ \\
\midrule
\multicolumn{3}{l}{\textit{Post-encounter (predicted)}} \\
$T_\mathrm{final}$ & $\sim 60\,\mathrm{K}$ & From adiabatic compression $\times 10$ \\
$n_\mathrm{final}$ & $\sim 10^5$--$10^6\cmcubed$ & Tidal compression \\
\bottomrule
\end{tabular}
\end{table}

%===============================================================================
\section{Three-Phase Simulation Workflow}
\label{sec:three_phase}
%===============================================================================

Based on the physics analysis above, we implement a three-phase workflow for
preparing stable initial conditions before the IMBH encounter.

\subsection{The Equilibrium Solution with $\gamma = 5/3$}

\textbf{Key insight}: Unlike the isothermal case, a $\gamma = 5/3$ polytrope
can achieve stable equilibrium for \textbf{any mass}. This resolves the fundamental
instability problem described in Section~\ref{sec:polytrope}.

For a Lane-Emden polytrope with $n = 1.5$ ($\gamma = 5/3$):
\begin{equation}
    P = K \rho^{5/3}, \quad \text{where} \quad K = \frac{k_B T}{\mu m_H \rho^{2/3}}
\end{equation}

The density profile is given by the Lane-Emden solution:
\begin{equation}
    \rho(r) = \rho_c \theta^{n}(\xi), \quad \xi = r/\alpha, \quad \alpha = \sqrt{\frac{(n+1)K\rho_c^{1/n-1}}{4\pi G}}
\end{equation}
where $\theta(\xi)$ satisfies the Lane-Emden equation with boundary conditions
$\theta(0) = 1$, $\theta'(0) = 0$.

\textbf{Advantages of $\gamma = 5/3$}:
\begin{itemize}
    \item Stable equilibrium for $M = 40\,\Msun$ (no critical mass limitation)
    \item Finite cloud radius ($\xi_1 = 3.654$ for $n = 1.5$)
    \item Adiabatic heating during compression naturally produces $T \sim 60\,\mathrm{K}$
\end{itemize}

\subsection{Phase 1: Relaxation (Particle Positioning)}

\textbf{Goal}: Create a particle distribution that matches the analytical density
profile with minimal numerical noise.

\begin{table}[h]
\centering
\caption{Phase 1: Relaxation Configuration ($\gamma = 5/3$ Polytrope)}
\begin{tabular}{lll}
\toprule
\textbf{Parameter} & \textbf{Value} & \textbf{Notes} \\
\midrule
EOS & $P = K\rho^{5/3}$ & Adiabatic (polytropic index $n = 1.5$) \\
Profile & Lane-Emden ($n = 1.5$) & $\rho = \rho_c \theta^{1.5}(\xi)$ \\
$M_\mathrm{cloud}$ & $40\,\Msun$ & Target mass \\
$T_\mathrm{initial}$ & $15\,\mathrm{K}$ & Sets polytropic constant $K$ \\
$n_\mathrm{center}$ & $10^3$--$10^4\cmcubed$ & Central number density \\
$R_\mathrm{cloud}$ & $0.5$--$1\pc$ & From Lane-Emden solution \\
\midrule
Self-gravity & \textbf{ON} & Required for true equilibrium \\
IMBH & OFF ($M_\mathrm{BH} = 0$) & No external perturbation \\
External envelope & ON (6 layers) & Provides pressure boundary \\
\midrule
Relaxation steps & 3000--5000 & Lane-Emden style \\
GLASS pre-relaxation & 1500--2000 & Uniform neighbor distribution \\
\bottomrule
\end{tabular}
\end{table}

\textbf{Relaxation method} (Lane-Emden style):
\begin{enumerate}
    \item Compute SPH forces (pressure + gravity)
    \item Subtract analytical equilibrium force: $\mathbf{a}_\mathrm{eq} = -\nabla P_\mathrm{eq}/\rho_\mathrm{eq}$
    \item Zero velocities each step (quasi-static relaxation)
    \item Update positions: $\Delta\mathbf{x} = \frac{1}{2}\mathbf{a}_\mathrm{net}\,\Delta t^2$
\end{enumerate}

\textbf{Output}: Relaxed particle distribution with density profile matching analytical solution.

\subsection{Phase 2: Hydrostatic Stability Test}

\textbf{Goal}: Verify that the relaxed cloud remains stable when evolved freely
with self-gravity (no IMBH).

\begin{table}[h]
\centering
\caption{Phase 2: Hydrostatic Test Configuration}
\begin{tabular}{lll}
\toprule
\textbf{Parameter} & \textbf{Value} & \textbf{Notes} \\
\midrule
Initial condition & Phase 1 output & Relaxed particle distribution \\
Self-gravity & \textbf{ON} & Test gravitational stability \\
IMBH & OFF ($M_\mathrm{BH} = 0$) & Isolate cloud dynamics \\
External envelope & ON & Maintains pressure boundary \\
Relaxation & OFF & Free evolution \\
\midrule
End time & $0.5$--$1.0\,\mathrm{Myr}$ & $\sim 0.5\,t_\mathrm{ff}$ \\
Output frequency & $0.05$--$0.1\,\mathrm{Myr}$ & Monitor evolution \\
\bottomrule
\end{tabular}
\end{table}

\textbf{Success criteria}:
\begin{itemize}
    \item Central density remains within $\pm 20\%$ of initial value
    \item Cloud radius does not change by more than $\pm 10\%$
    \item No runaway collapse or expansion
\end{itemize}

\textbf{Expected behavior}: Since the modified isothermal profile is not a
self-consistent equilibrium, the cloud will adjust slightly. This is acceptable
as long as the adjustment timescale is much longer than the IMBH encounter
timescale.

\subsection{Phase 3: IMBH Tidal Interaction}

\textbf{Goal}: Study the tidal disruption of the stable cloud by the IMBH.

\begin{table}[h]
\centering
\caption{Phase 3: IMBH Flyby Configuration}
\begin{tabular}{lll}
\toprule
\textbf{Parameter} & \textbf{Value} & \textbf{Notes} \\
\midrule
Initial condition & Phase 2 output & Verified stable cloud \\
$M_\mathrm{BH}$ & $5 \times 10^4\,\Msun$ & IMBH point mass \\
$b_\mathrm{impact}$ & $0$--$5\pc$ & Impact parameter \\
$v_\mathrm{rel}$ & $\sim 10\kms$ & Relative velocity \\
Self-gravity & ON & Cloud self-gravity \\
External envelope & ON & Pressure boundary \\
\midrule
Initial separation & $\sim 10\pc$ & Start before pericenter \\
End time & Until $\sim 10\pc$ post-pericenter & \\
\bottomrule
\end{tabular}
\end{table}

\textbf{Expected outcomes}:
\begin{itemize}
    \item Tidal stretching creates velocity gradient $\sim 100\kms/\pc$
    \item Dense clump forms near pericenter ($n \sim 10^{6.5}\cmcubed$)
    \item Cloud becomes gravitationally unbound ($M_\mathrm{vir}/M \sim 100$)
    \item Shock heating raises temperature to $\sim 60\,\mathrm{K}$
\end{itemize}

\subsection{Derived Cloud Parameters for $\gamma = 5/3$ Polytrope}

For a Lane-Emden polytrope with $n = 1.5$ ($\gamma = 5/3$):

\begin{table}[h]
\centering
\caption{Lane-Emden Polytrope Parameters ($\gamma = 5/3$, $M = 40\,\Msun$)}
\begin{tabular}{lll}
\toprule
\textbf{Parameter} & \textbf{Value} & \textbf{Derivation} \\
\midrule
EOS & $P = K\rho^{5/3}$ & Adiabatic polytrope \\
Polytropic index & $n = 1.5$ & $n = 1/(\gamma - 1)$ \\
Lane-Emden radius & $\xi_1 = 3.654$ & From LE solution \\
Mass parameter & $-\xi^2 \theta'|_{\xi_1} = 2.714$ & From LE solution \\
\midrule
Central density & $n_c = 10^3$--$10^4\cmcubed$ & Adjustable parameter \\
Cloud radius & $R = 0.5$--$1\pc$ & Depends on $n_c$ and $K$ \\
Initial temperature & $T_\mathrm{init} = 15$--$20\,\mathrm{K}$ & Sets $K$ \\
\midrule
Sound speed (center) & $c_s = 0.23\kms$ & At $T = 15\,\mathrm{K}$ \\
Sound crossing time & $t_\mathrm{sound} = 2$--$4\,\mathrm{Myr}$ & $R/c_s$ \\
Free-fall time & $t_\mathrm{ff} = 0.3$--$1\,\mathrm{Myr}$ & $\sqrt{3\pi/(32 G \bar{\rho})}$ \\
\midrule
\multicolumn{3}{l}{\textit{Post-encounter predictions (after compression $\times 10$)}} \\
Final temperature & $T_\mathrm{final} \sim 60\,\mathrm{K}$ & $T \propto \rho^{2/3}$ \\
Final density & $n_\mathrm{final} \sim 10^5$--$10^6\cmcubed$ & Tidal compression \\
\bottomrule
\end{tabular}
\end{table}

\subsection{Ghost Envelope Configuration}

The external pressure boundary is maintained by ghost particles arranged in
spherical shells outside the cloud:

\begin{table}[h]
\centering
\caption{Ghost Envelope Parameters}
\begin{tabular}{lll}
\toprule
\textbf{Parameter} & \textbf{Value} & \textbf{Notes} \\
\midrule
Number of layers & 6 & Sufficient for SPH kernel support \\
Layer spacing & $0.5\,h_\mathrm{edge}$ & Based on edge smoothing length \\
First layer & $R_\mathrm{cloud} + 0.5\,h_\mathrm{edge}$ & \\
Ghost density & $\rho_\mathrm{edge}$ & Matches cloud boundary \\
Ghost temperature & $T_\mathrm{cloud}$ & Isothermal \\
Ghost velocity & Fixed at 0 & Provides pressure, not mass flux \\
\bottomrule
\end{tabular}
\end{table}

\textbf{Physics}: Ghost particles provide the external pressure $P_\mathrm{ext}$
that confines the cloud. Without them, the cloud would expand into vacuum.

\subsection{Workflow Summary}

\begin{figure}[h]
\centering
\fbox{\parbox{0.95\textwidth}{
\textbf{Three-Phase Workflow ($\gamma = 5/3$ Polytrope + Cooling):}
\begin{enumerate}
    \item \textbf{Phase 1 (Relaxation)}: Create relaxed particle distribution
    \begin{itemize}
        \item EOS: $P = K\rho^{5/3}$ (adiabatic), \textbf{Cooling: OFF}
        \item Self-gravity ON, IMBH OFF
        \item Lane-Emden style relaxation (3000--5000 steps)
        \item Output: Particles positioned to match Lane-Emden $n = 1.5$ profile
    \end{itemize}
    \item \textbf{Phase 2 (Hydrostatic Test)}: Verify stability
    \begin{itemize}
        \item EOS: $\gamma = 5/3$, \textbf{Cooling: OFF}
        \item Self-gravity ON, IMBH OFF, relaxation OFF
        \item Run for $\sim 0.5$--$1\,t_\mathrm{ff}$
        \item Check: density profile stable (inherently stable for $\gamma = 5/3$)
    \end{itemize}
    \item \textbf{Phase 3 (IMBH Flyby)}: Study tidal disruption with cooling
    \begin{itemize}
        \item EOS: $\gamma = 5/3$, \textbf{Cooling: ON (KI2000)}
        \item Self-gravity ON, IMBH ON ($M_\mathrm{BH} \sim 5 \times 10^4\,\Msun$)
        \item Run through pericenter passage
        \item Observe: isothermal shock formation, dense clump ($n \to 10^{6.5}\cmcubed$)
    \end{itemize}
\end{enumerate}
}}
\end{figure}

%===============================================================================
\section{Recommended SPH Simulation Setup}
%===============================================================================

Based on the literature and first-principles analysis, we recommend:

\subsection{Baseline Configuration}

\begin{table}[h]
\centering
\caption{Recommended SPH Simulation Parameters}
\begin{tabular}{ll}
\toprule
\textbf{Parameter} & \textbf{Value} \\
\midrule
\multicolumn{2}{l}{\textit{Black Hole}} \\
$M_\mathrm{BH}$ & $10^5\,\Msun$ (Oka) or $5.7 \times 10^4\,\Msun$ (Ballone, minimum) \\
\midrule
\multicolumn{2}{l}{\textit{Cloud (Initial --- Pre-Encounter)}} \\
$M_\mathrm{cloud}$ & $40\,\Msun$ \\
$R_\mathrm{cloud}$ & $0.5$--$1\pc$ (larger than post-encounter) \\
$T_\mathrm{cloud}$ & $15$--$20\,\mathrm{K}$ (cold, pre-compression) \\
Density profile & Lane-Emden polytrope ($n = 1.5$) \\
Self-gravity & \textbf{ON} (required for stable equilibrium) \\
\midrule
\multicolumn{2}{l}{\textit{Orbit}} \\
Initial separation & $\sim 10\pc$ \\
Initial velocity & $\sim 10\kms$ toward BH \\
Impact parameter & $0$--$1\pc$ \\
\midrule
\multicolumn{2}{l}{\textit{Numerics}} \\
$N_\mathrm{particles}$ & $\geq 10^5$ \\
Equation of state & \textbf{Adiabatic ($\gamma = 5/3$)} \\
\bottomrule
\end{tabular}
\end{table}

\subsection{Validation Criteria}

The simulation is successful if:
\begin{enumerate}
    \item Post-encounter velocity width: $\Delta v \sim 100\kms$
    \item Line-of-sight velocity gradient: $\sim 70$--$130\kms/\pc$
    \item Compact structure forms: $R \lesssim 0.3\pc$
    \item Cloud is tidally disrupted (not self-bound)
\end{enumerate}

%===============================================================================
\section{Physical Interpretation}
%===============================================================================

\subsection{Scenario}

The HVCC CO--0.40--0.22 is interpreted as a molecular cloud that experienced a close encounter with an IMBH:

\begin{enumerate}
    \item \textbf{Pre-encounter}: A $\sim 40\,\Msun$ molecular cloud with $R \sim 0.2\pc$
    \item \textbf{Pericenter passage}: Strong tidal forces from $M_\mathrm{BH} \gtrsim 10^4\,\Msun$
    \item \textbf{Post-encounter}: Tidally stretched cloud with:
    \begin{itemize}
        \item Extreme velocity gradient ($\sim 100\kms$ over few pc)
        \item Compressed dense core ($n \sim 10^{6.5}\cmcubed$)
        \item Not gravitationally bound
    \end{itemize}
\end{enumerate}

\subsection{Why Not Other Explanations?}

\begin{itemize}
    \item \textbf{Supernova}: Would produce expanding shell, not linear velocity gradient
    \item \textbf{Stellar wind}: Insufficient energy for observed velocity width
    \item \textbf{Cloud-cloud collision}: Would not produce point-like radio source
    \item \textbf{Protostellar outflow}: Scale and morphology inconsistent
\end{itemize}

%===============================================================================
\section{Open Questions}
%===============================================================================

\begin{enumerate}
    \item What is the true $M_\mathrm{BH}$? (Range: $10^4$--$10^5\,\Msun$)
    \item What was the pre-encounter cloud structure?
    \item Is the cloud currently bound to the IMBH or escaping?
    \item How does adiabatic + shock heating combine to produce the observed 60 K?
    \item Can we reproduce the observed density enhancement ($10^4 \to 10^{6.5}\cmcubed$)?
\end{enumerate}

%===============================================================================
\section{Summary: Why $\gamma = 5/3$ Instead of Isothermal}
%===============================================================================

\begin{figure}[h]
\centering
\fbox{\parbox{0.95\textwidth}{
\textbf{Key Conclusions:}

\begin{enumerate}
    \item \textbf{Isothermal ($\gamma = 1$) fails}: For $M = 40\,\Msun$ and $T = 15\,\mathrm{K}$,
          the isothermal Bonnor-Ebert critical mass is only $\sim 1$--$2\,\Msun$.
          The cloud is 20--40$\times$ super-critical and must collapse.
          Numerical simulations confirmed this: timestep collapsed to $\sim 10^{-10}\,\mathrm{Myr}$.

    \item \textbf{Adiabatic ($\gamma = 5/3$) provides stability}: The Lane-Emden polytrope with $n = 1.5$
          has stable equilibrium for \textit{any} mass. The cloud can be set up in true
          hydrostatic equilibrium before the IMBH encounter.

    \item \textbf{Hybrid approach: $\gamma = 5/3$ + Cooling during flyby}:
          \begin{itemize}
              \item Phase 1--2: $\gamma = 5/3$, cooling OFF $\Rightarrow$ stable initial conditions
              \item Phase 3: $\gamma = 5/3$, cooling ON $\Rightarrow$ isothermal shocks form
          \end{itemize}

    \item \textbf{Isothermal shocks explain observations}:
          \begin{itemize}
              \item Adiabatic shocks: max compression $\rho_2/\rho_1 \leq 4$ (insufficient)
              \item Isothermal shocks: $\rho_2/\rho_1 = \mathcal{M}^2$ (can reach $n \sim 10^{6.5}\cmcubed$)
          \end{itemize}

    \item \textbf{Physical motivation}: Real molecular clouds have effective $\gamma > 1$
          due to turbulent pressure, magnetic fields, and finite cooling times.
          Cooling becomes important during strong compression (tidal encounter).
\end{enumerate}

\textbf{Bottom line}: The hybrid approach ($\gamma = 5/3$ for stability + cooling during flyby)
resolves the initial condition problem while naturally producing the observed dense clump
through cooling-triggered isothermal shocks.
}}
\end{figure}

%===============================================================================
\section{Critical Bonnor-Ebert Mass: Achieving 40 M$_\odot$ on the K\&I Curve}
\label{sec:be_critical_mass}
%===============================================================================

This section derives from first principles how to obtain a 40~$\Msun$ critical
Bonnor-Ebert sphere using the Koyama \& Inutsuka (2000) thermal equilibrium curve.

\subsection{The Bonnor-Ebert Critical Mass Formula}

For an isothermal self-gravitating sphere confined by external pressure $P_\mathrm{ext}$,
the maximum stable (critical) mass is:
\begin{equation}
    M_\mathrm{BE,crit} = 1.18 \frac{c_s^4}{\sqrt{G^3 P_\mathrm{ext}}}
    \label{eq:be_crit_mass}
\end{equation}

where the isothermal sound speed is:
\begin{equation}
    c_s = \sqrt{\frac{k_B T}{\mu m_H}}
    \label{eq:sound_speed}
\end{equation}

\subsubsection{Dependence on Mean Molecular Weight}

From Eq.~(\ref{eq:sound_speed}), we have $c_s \propto 1/\sqrt{\mu}$, therefore:
\begin{equation}
    M_\mathrm{BE,crit} \propto c_s^4 \propto \mu^{-2}
\end{equation}

\textbf{Critical insight}: Increasing $\mu$ from 1.27 (atomic) to 2.33 (molecular)
\textbf{decreases} the critical mass by a factor of:
\begin{equation}
    \frac{M_\mathrm{crit}(\mu=2.33)}{M_\mathrm{crit}(\mu=1.27)} = \left(\frac{1.27}{2.33}\right)^2 = 0.30
\end{equation}

This is why achieving 40~$\Msun$ is much easier with atomic gas than molecular gas.

\subsection{Scanning the K\&I 2000 Equilibrium Curve}

On the K\&I thermal equilibrium curve, $T_\mathrm{eq}(n)$ and $P_\mathrm{ext} = n k_B T_\mathrm{eq}$
are coupled. We compute $M_\mathrm{BE,crit}$ at various points along this curve.

\subsubsection{Derivation for Atomic Gas ($\mu = 1.27$)}

Using the K\&I 2000 data for $N_H = 10^{19}\,\mathrm{cm}^{-2}$:

\begin{table}[h]
\centering
\caption{K\&I 2000 Equilibrium and Critical BE Mass ($\mu = 1.27$, Atomic)}
\label{tab:ki2000_be_mass}
\begin{tabular}{cccccc}
\toprule
\textbf{$n$ [cm$^{-3}$]} & \textbf{$T_\mathrm{eq}$ [K]} & \textbf{$P/k_B$ [K cm$^{-3}$]} & \textbf{$c_s$ [km/s]} & \textbf{$M_\mathrm{BE,crit}$ [$\Msun$]} & \textbf{Phase} \\
\midrule
0.1 & 630 & 63 & 2.03 & $\sim 6 \times 10^4$ & WNM \\
1.0 & 60 & 60 & 0.63 & $\sim 6000$ & Unstable \\
10 & 20 & 200 & 0.36 & $\sim 200$ & Unstable \\
\textbf{50} & \textbf{9.4} & \textbf{470} & \textbf{0.25} & \textbf{$\sim 50$} & \textbf{CNM} \\
\textbf{70} & \textbf{8.75} & \textbf{612} & \textbf{0.24} & \textbf{$\sim 40$} & \textbf{CNM} \\
100 & 8.5 & 850 & 0.24 & $\sim 30$ & CNM \\
1000 & 7.3 & 7300 & 0.22 & $\sim 5$ & CNM \\
\bottomrule
\end{tabular}
\end{table}

\textbf{Result}: To achieve $M_\mathrm{BE,crit} \approx 40\,\Msun$ on the K\&I curve with
atomic gas ($\mu = 1.27$), we need:
\begin{itemize}
    \item Edge density: $n_\mathrm{edge} \approx 70\,\mathrm{cm}^{-3}$
    \item Temperature: $T \approx 9\,\mathrm{K}$
    \item External pressure: $P_\mathrm{ext}/k_B \approx 630\,\mathrm{K\,cm}^{-3}$
\end{itemize}

\subsubsection{Central Density for Critical BE Sphere}

For a critical Bonnor-Ebert sphere ($\xi_s = 6.45$), the density contrast is:
\begin{equation}
    \frac{\rho_\mathrm{center}}{\rho_\mathrm{edge}} = e^{\psi_s} \approx 14.04
\end{equation}

Therefore, if $n_\mathrm{edge} = 70\,\mathrm{cm}^{-3}$:
\begin{equation}
    n_\mathrm{center} = 14.04 \times 70 \approx 980\,\mathrm{cm}^{-3}
\end{equation}

\subsubsection{Cloud Radius}

The BE scale length is:
\begin{equation}
    r_0 = \frac{c_s}{\sqrt{4\pi G \rho_c}}
\end{equation}

For $c_s = 0.24\,\kms$, $n_c = 980\,\mathrm{cm}^{-3}$, $\mu = 1.27$:
\begin{align}
    \rho_c &= n_c \times \mu \times m_H = 980 \times 1.27 \times 1.67 \times 10^{-24}\,\mathrm{g/cm}^3 \\
           &= 2.08 \times 10^{-21}\,\mathrm{g/cm}^3 = 30.6\,\Msun/\mathrm{pc}^3
\end{align}

\begin{equation}
    r_0 = \frac{0.24}{\sqrt{4\pi \times 0.00430091 \times 30.6}} = \frac{0.24}{\sqrt{1.65}} = 0.19\,\mathrm{pc}
\end{equation}

The cloud radius:
\begin{equation}
    R_\mathrm{cloud} = \xi_s \times r_0 = 6.45 \times 0.19 \approx 1.2\,\mathrm{pc}
\end{equation}

\subsection{Options for a 40 M$_\odot$ Bonnor-Ebert Sphere}

\begin{table}[h]
\centering
\caption{Options for Achieving $M_\mathrm{BE} = 40\,\Msun$}
\label{tab:be_options}
\small
\begin{tabular}{lccccccl}
\toprule
\textbf{Option} & \textbf{$\mu$} & \textbf{$T$ [K]} & \textbf{$n_\mathrm{edge}$ [cm$^{-3}$]} & \textbf{$P/k_B$} & \textbf{$n_\mathrm{center}$ [cm$^{-3}$]} & \textbf{$R$ [pc]} & \textbf{Notes} \\
\midrule
\textcolor{blue}{\textbf{A}} & \textbf{1.27} & \textbf{9} & \textbf{70} & \textbf{630} & \textbf{980} & \textbf{1.2} & \textcolor{blue}{\textbf{On K\&I curve, atomic}} \\
B & 2.33 & 9 & 6 & 54 & 84 & 2.2 & Very low $P$, thermally unstable \\
C & 2.33 & 31 & 26 & 800 & 365 & 1.3 & Isothermal, NOT on K\&I curve \\
D & 2.33 & 31 & 250 & 7750 & 3500 & 0.65 & $M_\mathrm{crit} \approx 13\,\Msun$, not 40 \\
\bottomrule
\end{tabular}
\end{table}

\subsubsection{Option A: Atomic Gas on K\&I Curve (Recommended)}

\textbf{Physical justification}: This represents a Cold Neutral Medium (CNM) cloud that
is \textbf{not yet molecular}. The gas is primarily atomic hydrogen (HI) with $\mu = 1.27$.
After tidal compression by the IMBH, the increased density and shielding from UV radiation
can trigger H$_2$ formation, converting the cloud to molecular.

\begin{table}[h]
\centering
\caption{\textcolor{blue}{Recommended: Option A Parameters (Atomic CNM Cloud)}}
\begin{tabular}{lll}
\toprule
\textbf{Parameter} & \textbf{Value} & \textbf{Derivation} \\
\midrule
\multicolumn{3}{l}{\textit{Physical constants}} \\
Mean molecular weight & $\mu = 1.27$ & Atomic HI + 10\% He \\
Gravitational constant & $G = 0.00430091$ & Code units ($\Msun$, pc, km/s) \\
\midrule
\multicolumn{3}{l}{\textit{Edge (truncation) conditions from K\&I 2000}} \\
Edge density & $n_\mathrm{edge} = 70\,\mathrm{cm}^{-3}$ & K\&I equilibrium for $M_\mathrm{crit} = 40\,\Msun$ \\
Temperature & $T = 9\,\mathrm{K}$ & K\&I $T_\mathrm{eq}(n = 70)$ \\
External pressure & $P_\mathrm{ext}/k_B = 630\,\mathrm{K\,cm}^{-3}$ & $P = n_\mathrm{edge} \times T$ \\
Sound speed & $c_s = 0.24\,\kms$ & $\sqrt{k_B T / (\mu m_H)}$ \\
\midrule
\multicolumn{3}{l}{\textit{Central conditions (from BE profile, $\xi_s = 6.45$)}} \\
Density ratio & $\rho_c / \rho_\mathrm{edge} = 14.04$ & $e^{\psi_s}$ at critical truncation \\
Central density & $n_\mathrm{center} = 980\,\mathrm{cm}^{-3}$ & $14.04 \times n_\mathrm{edge}$ \\
Central T (K\&I) & $T_\mathrm{center} \approx 7.5\,\mathrm{K}$ & K\&I $T_\mathrm{eq}(n = 980)$ \\
\midrule
\multicolumn{3}{l}{\textit{Global cloud properties}} \\
Cloud mass & $M_\mathrm{cloud} = 40\,\Msun$ & Critical BE mass \\
Cloud radius & $R_\mathrm{cloud} = 1.2\,\mathrm{pc}$ & $\xi_s \times r_0$ \\
Scale length & $r_0 = 0.19\,\mathrm{pc}$ & $c_s / \sqrt{4\pi G \rho_c}$ \\
Stability & Critical ($M/M_\mathrm{crit} = 1$) & Marginally stable \\
\bottomrule
\end{tabular}
\end{table}

\subsubsection{Why Atomic Gas Works but Molecular Does Not}

The key issue is the strong dependence of $M_\mathrm{crit}$ on $\mu$:

\begin{equation}
    M_\mathrm{crit} \propto c_s^4 \propto T^2 / \mu^2
\end{equation}

\begin{table}[h]
\centering
\caption{Effect of Mean Molecular Weight on Critical Mass}
\begin{tabular}{lccc}
\toprule
\textbf{Gas Type} & \textbf{$\mu$} & \textbf{$c_s$ at $T = 10\,\mathrm{K}$ [km/s]} & \textbf{$M_\mathrm{crit}$ factor} \\
\midrule
Atomic (HI + He) & 1.27 & 0.26 & 1.0$\times$ \\
Molecular (H$_2$ + He) & 2.33 & 0.19 & 0.3$\times$ \\
\bottomrule
\end{tabular}
\end{table}

To achieve $M_\mathrm{crit} = 40\,\Msun$ with molecular gas ($\mu = 2.33$), you would need:
\begin{itemize}
    \item \textbf{Higher temperature}: $T \approx 30\,\mathrm{K}$ (not on K\&I curve at high density)
    \item \textbf{Lower pressure}: $P/k_B \sim 50$--$100\,\mathrm{K\,cm}^{-3}$ (very diffuse, unstable)
\end{itemize}

Neither option is physically realistic for a dense molecular cloud on the K\&I equilibrium curve.

\subsection{Simulation Strategy: Atomic $\to$ Molecular Transition}

The recommended approach uses atomic gas initial conditions that can become molecular
during the simulation:

\begin{enumerate}
    \item \textbf{Initial state}: Atomic CNM cloud at $T \approx 9\,\mathrm{K}$, $n_c \approx 1000\,\mathrm{cm}^{-3}$
    \item \textbf{IMBH compression}: Tidal forces compress cloud to $n \gtrsim 10^4\,\mathrm{cm}^{-3}$
    \item \textbf{H$_2$ formation}: At high density with sufficient shielding, H$_2$ forms
    \item \textbf{Final state}: Molecular cloud with $\mu \to 2.33$, observed as CO emission
\end{enumerate}

\textbf{Note}: The HVCC CO--0.40--0.22 is observed in CO emission, indicating molecular gas
in the \textbf{post-encounter} state. This is consistent with an initially atomic cloud
that became molecular during tidal compression.

\subsection{Code Unit Conversions for Option A}

For implementation in the SPH code (units: $\Msun$, pc, km/s):

\begin{table}[h]
\centering
\caption{Option A Parameters in Code Units}
\begin{tabular}{lcc}
\toprule
\textbf{Parameter} & \textbf{Physical Value} & \textbf{Code Units} \\
\midrule
Density-to-$n$ factor & -- & $32\,\mathrm{cm}^{-3}$ per $\Msun/\mathrm{pc}^3$ \\
$n_\mathrm{center} = 980\,\mathrm{cm}^{-3}$ & -- & $\rho_c = 30.6\,\Msun/\mathrm{pc}^3$ \\
$n_\mathrm{edge} = 70\,\mathrm{cm}^{-3}$ & -- & $\rho_\mathrm{edge} = 2.2\,\Msun/\mathrm{pc}^3$ \\
$c_s = 0.24\,\kms$ & -- & $c_s^2 = 0.058\,(\kms)^2$ \\
$P_\mathrm{ext}/k_B = 630\,\mathrm{K\,cm}^{-3}$ & -- & $P_\mathrm{ext} = \rho_\mathrm{edge} \times c_s^2$ \\
\bottomrule
\end{tabular}
\end{table}

%===============================================================================
\section{References}
%===============================================================================

\begin{enumerate}
    \item Oka, T., Tsujimoto, S., Iwata, Y., Nomura, M., \& Takekawa, S. (2017).
    ``Millimetre-wave Emission from an Intermediate-Mass Black Hole Candidate in the Milky Way.''
    \textit{Nature Astronomy}, 1, 709--712.
    \item Ballone, A., Mapelli, M., \& Pasquato, M. (2018).
    ``Tidal disruption of molecular clouds by intermediate-mass black holes.''
    \textit{Monthly Notices of the Royal Astronomical Society}, 480, 4684--4692.
    \item Oka, T., Mizuno, R., Miura, K., \& Takekawa, S. (2016).
    ``Signature of an Intermediate-mass Black Hole in the Central Molecular Zone of Our Galaxy.''
    \textit{Astrophysical Journal Letters}, 816, L7.
\end{enumerate}

\end{document}
